\documentclass[12pt,a4paper]{article}
\usepackage[utf8]{inputenc}
\usepackage[T1]{fontenc}
\usepackage[ngerman,latin]{babel}
\usepackage{amsmath,amssymb,amsthm}
\usepackage{geometry}
\usepackage{hyperref}
\usepackage{enumitem}
\usepackage{mathrsfs}
\usepackage{longtable}
\usepackage{booktabs}
\usepackage{microtype}
\usepackage{array}
\usepackage{multirow}
\geometry{margin=2.5cm}

\theoremstyle{plain}
\newtheorem{satz}{Satz}[section]
\newtheorem{theorem}{Theorem}[section]
\newtheorem{lemma}{Lemma}[section]
\newtheorem{hilfssatz}{Hilfssatz}[section]
\newtheorem{korollar}{Korollar}[section]

\theoremstyle{definition}
\newtheorem{definition}{Definition}[section]

\theoremstyle{remark}
\newtheorem{bemerkung}{Bemerkung}[section]
\newtheorem{beispiel}{Beispiel}[section]

\newcommand{\Q}{\mathfrak{Q}}
\newcommand{\fA}{\mathfrak{A}}
\newcommand{\fB}{\mathfrak{B}}
\newcommand{\fC}{\mathfrak{C}}

\title{Einzelband, Teil 9}
\author{Carl Friedrich Gau\ss}
\date{}

\begin{document}
\maketitle

\section*{Variae disquisitionum praecedentium applicationes}

\subsection*{Solutio aequationis indeterminatae $mxx+nyy=A$ per exclusiones.}

\subsubsection*{323.}

Omnes repraesentationes numeri dati $A$ per formam binariam $mxx+nyy$,
sive solutiones aequationis indeterminatae $mxx+nyy=A$, in Sectione V
methodo generali invenire docuimus, cuius brevitas quoque nihil desiderandum
relinquere videtur, si omnes valores expr.\ $\sqrt{-mn}$ secundum modulum $A$
ipsum, et per suos factores quadratos divisum, iam habentur; hic autem pro
eo casu, ubi $mn$ est positivus atque $n$ et $A$ esse positivos atque inter se
primos supponemus, ad hanc solutionem explicabimus directa multo
expeditiorem, si valores antea computare oportet. Supponemus autem,
numeros $m$, $n$ et $A$ esse positivos atque inter se primos, quum casus reliqui
ad hunc facile possint reduci.

Manifesto quoque sufficit, valores positivos ipsorum $x$, $y$ eruere, quum
reliqui inde per solam signorum mutationem deducantur.

Perspicuum est, $x$ ita comparatum esse debere, ut
$\frac{A-nyy}{m}=\frac{A}{m}-\frac{n}{m}yy$ fiat integer positivus, et quadratus evadat.

Conditio prima requirit, ut $x$ non sit maior quam
$\sqrt{\frac{A}{m}}$; secunda iam per se locum habet, quando $n=1$; quin requirit, ut
valor expr.\ $\frac{A-nyy}{m}$ sit residuum quadraticum ipsius $m$. Denotandoque
omnes valores diversos expr.\ $\sqrt{-n}(\text{mod.\ }m)$ per
$+r$, $-r$, $+r'$, $-r'$ etc., $x$ sub aliqua formarum
$nt+r$, $nt-r$, $nt+r'$ etc.\ contentus esse debebit.

Simplissimum itaque foret, omnes numeros harum formarum infra limitem
$\sqrt{\frac{A}{m}}$, quorum complexum per $Q$ exprimemus, pro $x$ substituere, eosque
solos retinere, pro quibus $\frac{A-nxx}{m}$ fit quadratum. Hoc tentamen, quantum
lubeat, contrahere, in art.\ sq.\ docebimus.

\subsubsection*{324.}

Methodus exclusionum, per quam hoc efficiemus, perinde ac in disqu.\ praec.
in eo consistit, ut plures numeros, etiam hic \emph{excludentes} vocandos,
accipiamus, pro quibusnam valoribus ipsius $x$ valor ipsius
$\frac{A-nxx}{m}$ fiat non-residuum, talesque ad lubitum ex $Q$ eiiciamus.

Per rationem art.\ 321, 322 iis quae in art.\ 321 sqq.\ exposuimus, omnino
analoga apparet, tales tantum excludentes adhibendos esse, qui sint numeri
primi aut numerorum primorum potestates, et pro excludente posterioris
generis ea tantum ipsius non-residua a valoribus ipsius $x$ arcenda, quae sint
residua omnium potestatum inferiorum eiusdem numeri primi, siquidem
exclusio cum his iam est instituta.

Si itaque excludens $E=p^\lambda$, ubi $p$ est numerus primus ipsum $m$ non
metiens, supponamusque brevitatis caussa duos casus, in quibus $n$ per $p$
est divisibilis ac non divisibilis, simul complectimur (includendo etiam casum
$\lambda=1$), esse summam potestatem eiusdem numeri primi per quam $n$ sit
divisibilis. Sint $a$, $b$, $c$ etc.\ non-residua quadratica ipsius $E$ (omnia,
quando $n$ per $p$ non est divisibilis; ea tantum, quae simul sunt residua
potestatum inferiorum, quando $n$ per $p$ est divisibilis).\ Computentur radices
congruentiarum $mz\equiv A-na$, $mz\equiv A-nb$, $mz\equiv A-nc$ etc.\ (mod.\ $E$),
sive ea quae necessaria sive ea quae sufficiunt; quae sint $\alpha$, $\beta$, $\gamma$ etc.\ patetque
facile, si quis valor ipsius $x$ producat $xx\equiv\alpha$ (mod.\ $E$) sive non-residuum
ipsius $E$, similiterque de numeris reliquis $\beta$, $\gamma$ etc.\ aeque facile vice versa
perspicitur, si quis valor ipsius $x$ producat $xx\equiv\alpha$ (mod.\ $E$), pro eodem
fieri $\frac{A-nxx}{m}\equiv a$ (mod.\ $E$), adeoque omnes valores ipsius $x$, pro quibus
$xx$ nulli numerorum $a$, $b$, $c$ etc.\ sec.\ mod.\ $E$ congruus sit, tales valores
ipsius $x$ producere, qui nulli numerorum $a$, $\beta$, $\gamma$ etc.\ sec.\ mod.\ $E$ sint
congrui.

Eligantur iam e numeris $a$, $\beta$, $\gamma$ etc.\ omnia residua quadratica ipsius $E$,
quae sint $h$, $h'$, $h''$ etc., computentur valores expressionum $\sqrt{h}$,
$\sqrt{h'}$, $\sqrt{h''}$ etc.\ (mod.\ $E$), ponamusque hinc prodire $+a$, $+a'$, $+a''$
etc.\ His ita factis manifestum est, omnes numeros formarum
$Et+a$, $Et+b$, $Et+c$ etc., $Et+h$, $Et+h'$, $Et+h''$ etc.\ ex $Q$ tuto eiici
posse, nullique valori ipsius $x$ post hanc exclusionem remanenti valorem
ipsius $\frac{A-nxx}{m}$ sub formis $Et+a$, $Et+b$, $Et+c$ etc.\ contentum respondere
posse.

Ceterum manifestum est, tales valores ipsius $x$ iam per se e nullo valore
ipsius $x$ prodire posse, quando inter numeros $a$, $\beta$, $\gamma$ etc.\ nulla residua qu.\ ipsius $E$ inveniantur, adeoque in hoc casu numerum $E$ tamquam excludentem
applicari non posse.

Huiusmodi excludentes, quot placet, adhiberi, atque sie numeri in $Q$ ad
lubitum diminui possunt.

Videamus iam, annon etiam numeros primos ipsum $m$ metientes, taliumve
numerorum potestates tamquam excludentes adhibere liceat. Sit $B$ valor
expr.\ $-\frac{A}{m}$ (mod.\ $m$), patetque, $\frac{A-nxx}{m}$ semper ipsi $B$ secundum
mod.\ $m$ congruum fieri, quicunque valor pro $x$ accipiatur, adeoque ad
possibilitatem aequ.\ prop.\ necessario requiri, ut $B$ sit residuum
quadraticum ipsius $m$.

Designante itaque $p$ quemcunque primum imparem ipsius $m$ qui per hyp.\ ipsos $n$ et $A$, adeoque etiam ipsum $B$ non metietur, pro valore quocunque
ipsius $x$ erit $\frac{A-nxx}{m}$ residuum ipsius $p$ adeoque etiam cuiuscunque
potestatis ipsius $p$; quamobrem $p$ ipsiusque potestates nequeunt
excludentium loco haberi, quando $m$ per $p$ est divisibilis.

Prorsus simili ratione, si $m$ per 4 est divisibilis, ad aequ.\ prop.\ possibilitatem necessario requiritur, ut sit $B\equiv 1$ (mod.\ 4), unde etiam
$\frac{A-nxx}{m}$ pro valore quocunque ipsius $x$ fiet $\equiv B\equiv 1$ (mod.\ 4), et proin
binarii potestates ad exclusionem non idoneae, quando $m$ per 4 est divisibilis.

Quando autem $m$ per 8 est divisibilis, ad aequ.\ prop.\ possibilitatem
necessario requiritur, ut sit $B\equiv 1$ (mod.\ 8), unde etiam
$\frac{A-nxx}{m}$ pro valore quocunque ipsius $x$ fiet $\equiv B\equiv 1$ (mod.\ 8), et
proin binarii potestates ad exclusionem non idoneae.

Quando autem $m$ per 4 neque vero per 8 est divisibilis, ex simili ratione
esse debebit $B\equiv 1$ (mod.\ 4), adeoque valor expr.\ $-\frac{A}{m}$ (mod.\ 8) vel $=1$
vel $=5$, designetur per $C$; pro $x$ impari, $n$ pari, $\frac{A-nxx}{m}$ fit $\equiv C$;
pro $x$ pari, $\frac{A-nxx}{m}$ fit $\equiv C+4$ (mod.\ 8). Nullo negotio perspicitur,
valores impares reiiciendos esse, quando $C=5$; omnes pares, quando $C=1$;
omnes pares, quando $C=5$.

Denique quando $m$ per 2 neque vero per 4 est divisibilis, sit ut ante $C$
valor expr.\ $-\frac{A}{m}$ (mod.\ 8), qui erit 1, 3, 5 vel 7; atque $D$ valor huius
$-\frac{n}{m}$ (mod.\ 4), qui erit 1 vel 3. Facile hinc colligitur, omnes valores impares
ipsius $x$, quando $C=1$ et $D=3$, aut $C=7$ et $D=3$, omnes valores pares,
quando $C=3$ et $D=1$, aut $C=5$, aut $C=3$ et $D=3$, aut $C=7$ et $D=1$,
fieri $\frac{A-nxx}{m}\equiv 3$, 5 vel 7 (mod.\ 8), adeoque reiiciendos esse; in casibus
reliquis autem, puta quando $C=1$ et $D=1$, aut $C=7$ et $D=1$, fiet
$\frac{A-nxx}{m}\equiv 1$ (mod.\ 8) pro omnibus valoribus ipsius $x$, adeoque pro $x$
pari $\equiv C$, pro impari $\equiv C$ (mod.\ 8); quum itaque valor ipsius
$\frac{A-nxx}{m}$ semper fiat $\equiv 1$ (mod.\ 8), sive $x$ accipiatur par sive impar,
unde liquet, in his casibus aequationem prop.\ solutionem omnino non admittere.

Ceterum quum prorsus simili modo, ut hic valorem ipsius $x$ per exclusiones
invenire docuimus, etiam, mutatis mutandis, valorem ipsius $y$ elicere
possimus, methodum exclusionis ad problematis propositi solutionem duobus
semper modis applicare licebit (nisi $m=n=1$, ubi coincidunt), e quibus
is plerumque est praeferendus, pro quo $Q$ terminorum multitudinem minorem
continet, quod facile a priori aestimari poterit.

Denique vix necesse erit observare, si post aliquot exclusiones omnes
numeri ex $Q$ abierint, hoc ut certum indicium impossibilitatis aequationis
propositae esse considerandum.

\subsubsection*{325.}

\textbf{Ex.}\ Proposita sit aequatio $455xx+3yy=10857362$, quam duplici
modo solvemus, primo investigando valores ipsius $x$, dein valores ipsius $y$.
Limes illorum in hoc casu est $\sqrt{\frac{10857362}{455}}$, qui cadit inter 1902 et 1903;
valor expr.\ $\sqrt{-3}$ (mod.\ 455) est $\pm 82$, $\pm 152$, $\pm 173$, $\pm 212$; atque
valores expr.\ $\sqrt{36191201-455y}$ (mod.\ 354) sunt $\pm 173$, $\pm 212$.

Hinc $Q$ constat e 33 numeris sequentibus:

82, 152, 173, 212, 243, 282, 303, 373, 537, 607, 628, 667, 698, 737,
758, 828, 992, 1062, 1083, 1122, 1153, 1192, 1213, 1283, 1447, 1517,
1538, 1577, 1608, 1647, 1668, 1738, 1902.

Numerus 3 in hoc casu ad exclusionem adhiberi nequit, quia ipsum $m$
metitur.

Pro excludente 4 habemus $a=2$, $b=3$, unde $\alpha=2$, $\beta=3$; $y'y$ (mod.\ 4)
fit 0, 1; atque valores expr.\ $\sqrt{0}$, $\sqrt{1}$ (mod.\ 4) hos $0$, $1$; hinc sequitur,
omnes numeros formarum $4t$ et $4t+2$, i.\ e.\ omnes pares ex $Q$ eiiciendos
esse; designentur $Q'$ numeri formarum $4t+1$ et $4t+3$ (sedecim) reliqui
per hos:

$Q'$: 173, 373, 537, 607, 628, 667, 698, 737, 758, 828, 992, 1062, 1083,
1122, 1153, 1192, 1213, 1283, 1447, 1517, 1538, 1577, 1608, 1647, 1668,
1738, 1902.

Pro $E=5$, qui etiam ipsum $n$ metitur, habemus $a=2$, $b=3$, unde
$\alpha=8$, $\beta=9$; radices congruentiarum $mz\equiv A-2n$ et $mz\equiv A-3n$
(mod.\ 25) has 9 et 24, quae ambae sunt residua ipsius 25, valoresque
expressionum $\sqrt{9}$ et $\sqrt{24}$ (mod.\ 25) hos $\pm 3$, $\pm 7$; eiectis ex $Q'$ omnibus
numeris formarum $25t+7$, $25t+3$, $25t+17$, $25t+23$ remanent hi decem
($Q''$):

173, 373, 537, 667, 737, 1083, 1213, 1283, 1517, 1577.

Pro $E=7$, habemus congruentiarum $mz\equiv A-an$, $mz\equiv A-bn$,
$mz\equiv A-\beta n$ (mod.\ 49) radices 32, 39, 18, quae omnes sunt residua
ipsius 49, atque valores expr.\ $\sqrt{32}$, $\sqrt{39}$, $\sqrt{18}$ (mod.\ 49) hos $\pm 9$, $\pm 23$,
$\pm 19$; eiectis ex $Q''$ numeris formarum $49t+9$, $49t+19$, $49t+23$
remanent hi quinque ($Q'''$): 537, 737, 1083, 1213, 1517.

Pro $E=8$, habemus $a=5$, qui est non-residuum ipsius 8; unde $a=5$,
$b=10$; residua ipsius 8 puta $1$, $4$; unde $\alpha=5$, $\beta=0$; hinc deducitur,
ex $Q'''$ reiiciendos esse numeros formarum $11t$, $11t+1$, $11t+5$, quo facto
remanent 537, 1083, 1213. Quos tentando prodeunt pro $V$ resp.\ valores
21961, 16129, 14161, e quibus secundus ac tertius soli sunt quadrata. Quare
aequ.\ prop.\ duas solutiones per valores positivos ipsorum $x$, $y$ admittit:
$x=1083$, $y=127$, et $x=1213$, $y=119$.

Secundo.\ Si alteram eiusdem aequationis incognitam per exclusiones
indagare placet, ponatur haec sub formam $455yy+3xx=10857362$,
commutando $x$ cum $y$, ut omnia signa artt.\ 323, 324 retinere liceat. Limes
valorum ipsius $x$ hic cadit inter 154 et 155; valor expr.\ $-\frac{A}{m}$ (mod.\ $m$)
est 1; valores huius $\sqrt{-1}$ (mod.\ 3) sunt $+1$ et $-1$, i.\ e.\ $3t-1$,
$3t+1$. Quare $Q$ omnes numeros formarum $3t+1$ et $3t+2$, i.\ e.\ omnes
per 3 non divisibiles usque ad 154 incl.\ continet; quorum multitudo est 103;
applicando autem praecepta supra data invenitur reiiciendos esse numeros
formarum $4t+2$ sive omnes pares; pro excl.\ 3; 4; 9; 11; 17; 19; 23:
$9t+7$; $11t$, $11t+3$; $17t+3$, $17t+4$, $17t+5$, $17t+7$; $19t+2$,
$19t+3$, $19t+8$, $19t+9$; $23t$, $23t+1$, $23t+5$, $23t+7$, $23t+9$,
$23t+10$.

His deletis superstites inveniuntur 119, 127, qui duo soli ipsi $V$ valorem
quadratum conciliant, easdemque solutiones suggerunt, ad quas supra
pervenimus.

\subsubsection*{326.}

Methodus praecedens iam per se tam expedita est ut vix quidquam
optandum relinquat; attamen per multifaria artificia magnopere adhuc
contrahi potest, e quibus hic pauca tantum attingere licet.

Restringemus itaque disquisitionem ad eum casum, ubi excludens est
numerus primus impar ipsum $A$ non metiens, sive talis primi potestas,
praesertim quoniam casus reliqui vel ad hunc reduci vel methodo analoga
tractari possunt.

Supponendo $E=p$ esse numerum primum ipsos $m$, $n$ non
metientem, atque valores expr.\ $-\frac{A}{m}$, $-\frac{A}{m}+\frac{n}{m}$, $-\frac{A}{m}+\frac{2n}{m}$ etc.\ (mod.\ $p$)
resp.\ numeri $k$, $\mathcal{E}$, $\mathfrak{E}$ etc., numeri $a$, $b$, $\gamma$ etc.\ inveniuntur per congruentias
$\alpha\equiv k+\mathfrak{A}$, $\beta\equiv k+\mathfrak{B}$, $\gamma\equiv k+\mathfrak{C}$ etc.\ (mod.\ $p$).\ Numeri $\mathfrak{A}$,
$\mathfrak{B}$, $\mathfrak{C}$ etc.\ autem per artificium ei prorsus simile, quo in art.\ 322 usi sumus,
sine congruentiarum computatione erui possunt, et vel cum omnibus
non-residuis, vel cum omnibus residuis ipsius $p$ (praeter 0) convenient, prout
valor expr.\ $-\frac{n}{m}$, sive (quod hie eodem redit) numerus $-4mn$ est residuum
vel non-residuum ipsius $p$.

Ita in ex.\ II art.\ praec.\ pro $E=11$ hinc numeri $\mathfrak{A}$, $\mathfrak{B}$ etc.\ erunt
$1$, $2$, $4$, $8$, $9$, $13$, $15$, $16$ adeoque numeri $a$, $b$ etc.\ $8$, $9$, $11$, $15$, $16$,
$3$, $5$, $6$; ex his residua sunt $8$, $9$, $11$, $15$, $16$, unde $+a$, $+a'$, $+a''$ etc.\ fiunt $\pm 4$, $\pm 5$.

Ceterum observamus adhuc, multitudinem numerorum $a$, $a'$, $a''$ etc.
semper esse $\frac{1}{2}(p-1)$, quando uterque numerus $k$ et $-mn$ est vel uterque
residuum vel uterque non-residuum ipsius $p$; $\frac{1}{2}(p+1)$, quando prior NR.,
posterior R.; $\frac{1}{2}(p-3)$, quando prior R., posterior NR.\ Demonstrationem
huius theorematis, ne nimis prolixi fiamus, supprimere debemus.

Quod autem secundo eos casus attinet, ubi $E$ est numerus primus ipsum
$n$ metiens, aut potestas numeri primi (imparis) ipsum $n$ metientis seu non
metientis, hi adhuc expeditius tractari possunt.

Omnes hos casus simul complectemur, omnibusque art.\ 324 signis
retentis ponemus $n=\mu p^\lambda$, ita ut $n'$ per $p$ non sit divisibilis. Numeri $a$,
$b$, $c$ etc.\ erunt producta numeri $p^{\lambda-1}$ vel in omnes numeros ipso $p$
minores (praeter 0), vel in omnia non-residua ipsius $p$ infra $p$, prout
$\mu$ est par vel impar; vel in omnia residua, quando $-\frac{n'}{m}$ est par vel impar;
vel in omnes numeros infra $p$ (praeter 0), quando $\mu$ impar atque $-n'm$ R$p$;
vel in omnia non-residua ipsius $p$ infra hunc limitem, quando $\mu$ impar
atque $-n'm$ N$p$.

Sit $k$ valor expr.\ $-\frac{A}{m}$ (mod.\ $p^{\lambda+\mu}$), eritque per $p$ non divisibilis,
quia eadem proprietas in $A$ supponitur; porro patet, omnes $\alpha$, $\beta$, $\gamma$
etc.\ ipsi $k$ sec.\ mod.\ $p$ congruos fieri, adeoque $p^\lambda$ nihil ex $Q$ excludere,
si $k$ N$p$; si vero $k$ R$p$ adeoque etiam $k$ R$p^{\lambda+\mu}$, sit $r$ valor expr.\ $\sqrt{k}$
(mod.\ $p^{\lambda+\mu}$), qui per $p$ non erit divisibilis, atque $e$ valor huius
$\frac{n'}{2mr}$ (mod.\ $p$), eritque $\alpha\equiv rr+2erp^\lambda$ (mod.\ $p^{\lambda+\mu}$), unde facile
colligitur, $\alpha$ esse residuum ipsius $p^{\lambda+\mu}$, atque valores expr.
$\sqrt{\alpha}$ (mod.\ $p^{\lambda+\mu}$) fieri $\pm(r+ep^\lambda+up^{\lambda+\mu})$; hinc omnes $a$, $a'$, $a''$
etc.\ exprimentur per $r+up^{\lambda+\mu}$. Denique nullo negotio hinc concluditur,
numeros $h$, $h'$, $h''$ etc.\ oriri ex additione numeri $r$ cum productis
numeri $p^\lambda$ vel in omnes numeros infra $p$ (praeter 0) puta quando $\mu$ impar,
$k$ R$p$; vel in omnia non-residua ipsius $p$ infra hunc limitem, quando
$\mu$ impar atque $-n'm$ N$p$; vel in omnia residua (praeter 0), quando
$\mu$ impar atque $-n'm$ R$p$, sive, quod hie eodem redit, quando $-n'm$ R$p$
sive $\mu$ par.

Ceterum simulac pro singulis excludentibus, quos applicare placent,
numeri $a$, $h'$, $h''$ etc.\ sunt eruti, quos ex $Q$ excludere oportet, exclusionem
ipsam etiam per operationes mechanicas perficere licebit, quales quisque
harum rerum peritus facile proprio marte excogitare poterit, si operae
pretium esse videbitur.

Tandem observare debemus, quamvis aequationem $axx+2bxy+cyy=M$,
in qua $bb-ac$ negativus sit, facile ad eam formam, quam in praecc.
consideravimus, reduci posse. Designando enim divisorem communem
maximum numerorum $a$, $b$ per $m$, et ponendo $a=md$, $b=mb'$,
$ac-bb'=de-mb'b'=n$, manifesto aequ.\ illa aequivalet huic
$da'x+b'y^2+nyy=dM$, quae per praecepta supra tradita solvi poterit.\ Ex
huius autem solutionibus eae tantum erunt retinendae, in quibus $x'-b'y$ per
$d$ fit divisibilis, sive unde $x$ valores integres nanciscitur.

\subsection*{Alia methodus congruentiam $xx\equiv A$ solvendi pro eo casu,
uhi $A$ est negativus.}

\subsubsection*{327.}

Quemadmodum solutio directa aequationis $axx+2bxy+cyy=M$ in
Sect.\ V contenta valores expr.\ $\sqrt{bb-ac}$ (mod.\ $M$) notos supponit: ita
vice versa pro eo casu, ubi $bb-ac$ est negativus, solutio indirecta in
praecc.\ exposita methodum expeditissimam subministrat, illos valores
eruendi, quae, praesertim pro valore permagno ipsius $M$, methodo art.\ 322
sqq.\ longo est praeferenda.

Supponemus autem, $M$ esse numerum primum, aut saltem ipsius
factores, si compositus esset, adhuc incognitos esse; enim si constaret
$M=pM'$, numerum primum $p$ ipsum $M$ metiri, atque esse $M'$ factorem
qui $p$ non amplius implicet, ita ut $M$ longe commodius foret, valores expr.\ $\sqrt{-D}$ pro modulis $p^\lambda$ et $M'$ sigillatim explorare (priores ex valoribus
secundum modulum $p$, art.\ 101), valoresque sec.\ mod.\ $M$ ex horum
combinatione deducere (art.\ 105).

Quaerendi sint itaque omnes valores expr.\ $\sqrt{-D}$ (mod.\ $M$), ubi $D$ et
$M$ positivi supponuntur (art.\ 147 sqq.), alioquin enim a priori constaret,
nullos numeros expressioni propositae satisfacere posse. Sint valores quaesiti,
e quibus bini semper oppositi erunt, $\pm r$, $\pm r'$, $\pm r''$ etc., atque $D+rr=Mh$,
$D+r'r'=Mh'$, $D+r''r''=Mh''$ etc.; porro designentur classes, ad quas
formae $(M,-r,h)$, $(M,r,h)$, $(M,-r',h')$, $(M,r',h')$, $(M,-r'',h'')$
etc.\ pertinent, resp.\ per $\mathfrak{C}$, $\mathfrak{C}'$, $\mathfrak{C}''$ etc., ipsarumque complexus per $\mathfrak{G}$.

Hae classes quidem, generaliter loquendo, tamquam incognitae sunt
spectandae; attamen perspicuum est, primo, omnes esse positivas atque
proprie primitivas, secundo, omnes ad idem genus pertinere, cuius character
ex indole numeri $M$, i.\ e.\ ex ipsius relationibus ad singulos divisores primos
ipsius $D$ (insuperque ad 4 aut 8, quando hae sunt necessariae) facile
cognosci possit (art.\ 230).

Quum suppositum sit, $M$ contineri sub forma divisorum ipsius
$xx+D$, a priori certi esse possumus, huic characteri necessario satisfieri,
nequeat; quum itaque hoc genus sit notum, omnes classes in ipso contentae
erui poterunt, quae sint $C$, $C'$, $C''$ etc., atque ipsarum complexus $G$.
Patet igitur, singulas classes $\mathfrak{C}$, $\mathfrak{C}'$, $\mathfrak{C}''$ etc.\ cum aliqua classe in
$G$ identicas esse debere; fieri potest quoque, ut plures classes in $\mathfrak{G}$
inter se, adeoque cum eadem in $G$ identicae sint, et quando $G$ unicam
classem continet, certo omnes in $\mathfrak{G}$ cum hac convenient. Quare si e
classibus $C$, $C'$, $C''$ etc.\ formae (simplissimae) $f$, $f'$, $f''$ etc.\ in
$\mathfrak{G}$ una forma inter has reperietur.

Iam si $axx+2bxy+cyy$ est forma in classe $\mathfrak{C}$ contenta, dabuntur
duae repraesentationes numeri $M$ per ipsam ad valorem $r$ pertinentes, et si
una est $x=m$, $y=n$, altera erit $x=-m$, $y=-n$; unicus casus
excipi debet ubi $D=1$, in quo quatuor repraesentationes dabuntur (v.\ art.\ 180).

Ex his colligitur, si omnes repraesentationes numeri $M$ per singulas
formas $f$, $f'$, $f''$ etc.\ investigentur (per methodum indirectam in praecc.
traditam), atque hinc valores expr.\ $\sqrt{-D}$ (mod.\ $M$) ad quos singulae
pertinent deducantur (art.\ 154 sqq.), omnes valores huius expressionis inde
obtineri, et quidem singulos bis, aut, si $D=1$, quater.\ Q.\ E.\ F.\ Si
quae formae inter $f$, $f'$, $f''$ etc.\ reperiuntur, per quas $M$ repraesentari
nequit, hoc est indicium, ipsas ad nullam classem in $\mathfrak{G}$ pertinere, adeoque
neglegendas esse: si vero $M$ per nullam illarum formarum repraesentari
potest, necessario $-D$ debebit esse non-residuum quadraticum ipsius $M$.

Circa has operationes teneantur adhuc observationes sequentes.

I.\ Repraesentationes numeri $M$ per formas $f$, $f'$ etc., quas hic
adhibemus, subintelliguntur esse tales, in quibus indeterminatarum valores
inter se primi sunt; si quae aliae se offerunt, in quibus hi valores divisorem
communem $\mu$ habent (quod tunc tantummodo accidere potest, ubi $\mu^2$ metitur
ipsum $M$, certoque accidet, quando $-D$ R$M$): hae ad institutum praesens
omnino negligi debent, etsi alio respectu utiles esse possint.

Ceteris paribus labor manifesto eo facilior erit, quo minor est multitudo
classium $f$, $f'$, $f''$ etc., adeoque brevissimus, quando $D$ est unus e 65
numeris in art.\ 303 traditis, pro quibus in singulis generibus unica tantum
classis datur.

II.\ Quum binae semper huiusmodi repraesentationes $x=m$, $y=n$;
$x=-m$, $y=-n$ ad eundem valorem pertineant, perspicuum est,
sufficere, si eae tantummodo repraesentationes considerentur, in quibus $y$
positivus.\ Tales itaque repraesentationes diversae semper valoribus diversis
expr.\ $\sqrt{-D}$ (mod.\ $M$) respondent, unde multitudo omnium valorum
diversorum multitudini omnium talium repraesentationum prodeuntium
aequalis erit (excepto casu $D=1$).

III.\ Manifesto etiam vicissim, si de duabus classis eiusdem formis,
puta $f$ et $f'$, altera e reliquis omnibus signis excepto eo, quod ad $x$
pertinet, maneat, repraesentationes per $f'$ oriri e repraesentationibus per $f$,
si in his valor ipsius $x$ signo mutetur; ex his vero easdem valores
expressionis $\sqrt{-D}$ (mod.\ $M$) prodire, quales ex illis manant, facile
perspicietur.\ Quamobrem inter formas $f$, $f'$, $f''$ etc.\ nonnisi semissem
earum considerare opus est.

IV.\ Quum pro valore dato ipsius $x$ valor ipsius $y$ per congruentiam
$2bxy\equiv M-axx-D$ (mod.\ $m$) vel omnino non detur, vel duos valores
inter se oppositos habeat, perspicuum est, in repraesentationibus numeri $M$
per formam $(M,-r,h)$, in quibus $y$ est positivus, ipsum $y$ necessario
infra $M$ contentum esse; hunc vero $2b$ terminorum limitem adhuc multo
magis restringere licebit, si pro $y$ accipiatur residuum minimum positivum
congruentiae $y\equiv\frac{M-axx-D}{2bx}$ (mod.\ $m$), ita ut $x$ infra terminum
$\sqrt{\frac{M}{a}}$ atque $y$ infra terminum $2b$ sit contentus.\ In hac suppositione
autem valores ipsius $x$, $y$ unius repraesentationis necessario a valore ipsius
$x$ alterius diversi erunt, ita ut, excepto casu $D=1$, repraesentationes, e
quibus valores diversi $\sqrt{-D}$ (mod.\ $M$) oriuntur, etiam inter se diversae
esse debeant.

V.\ Denique observamus, quum ad solutionem propositam sufficiat,
si pro unaquaque forma $f$ etc.\ unica tantum repraesentatio habeatur,
quam ex qua data elici possit, manifesto valorem ipsius $x$ aliquantisper infra
$\sqrt{\frac{M}{a}}$ descendere oportere, si pro valore ipsius $y$ valorem minimum
proprium adoptare velimus.\ Itaque si sub forma $f$ haec est
$(a,b,-a')$, faciendum erit, ut fiat $axx+2bxy-a'yy=M$, seu
$(ax+by)^2+Dyy=aM$, adeoque manifesto necesse est, ut $yy$ ipsum $aM$
non excedat, unde elicitur $y\leq\sqrt{\frac{aM}{D}}$, atque pro $x$ valorem minimum
proprium, si quidem fieri potest, eligere oportet.\ His limitibus observatis,
si valor ipsius $x$ infra $0$ descendit, repraesentatio ipsi $f$ omnino competit
non potest, adeoque eiusmodi forma ad solutionem nostram omnino non
pertinet.

\subsubsection*{328.}

\textbf{Ex.}\ Sint quaerendi omnes valores expr.\ $\sqrt{-1365}$ (mod.\ $5428681=M$).
Valor expr.\ $-\frac{D}{m}$ hic est 5428681.\ Valor characteris generis, ad quem
classes quaesitae pertinent, e relationibus numeri $M$ ad 3, 5, 7, 13 invenitur
$=1$, $4$; $1$, $4$; $1$, $4$; $1$, $4$, ita ut hic genus per 3, 5, 7, 13
R$M$; $M$ R$3$, $M$ R$5$, $M$ R$7$, $M$ R$13$ determinari possit.

Huius generis classes sunt quatuor, e quibus eligendae formae simplissimae
$f=(1,1,341)$, $f'=(9,5,151)$, $f''=(15,5,91)$, $f'''=(15,-5,91)$.

Ex repraesentationibus per $f$ oriuntur valores expr.\ $\sqrt{-D}$ (mod.\ $M$),
qui sunt $\pm 310439$, $\pm 2340424$, $\pm 3224413$.

Ex repraesentationibus per $f'$ prodierunt valores $\pm 1176000$,
$\pm 2327472$, $\pm 3192457$, $\pm 467334$.

Ex repraesentationibus per $f''$ et $f'''$ (quarum altera tantummodo opus
est) prodeunt valores $\pm 1748204$, $\pm 1702845$, $\pm 2955548$, $\pm 500759$.

Omnes duodecim valores diversi sunt, demtis signis.\ Combinando eos
cum valoribus eodem modo erutis pro modulis 4 et 8 (quum hic $D$ sit
$\equiv 3$ (mod.\ 4) et $M\equiv 1$ (mod.\ 8)), obtinentur omnes sedecim valores
quaesiti $\pm r$, $\pm r'$ etc.

\subsubsection*{329.}

Quum methodus art.\ praec.\ nitatur resolutione formae $(a,b,c)$ in
 $(a,b,-a')$, eadem fere laborat difficultate, qua methodus directa in
Sectione V tradita premitur.\ Adeoque haud iniuria desiderari posset alia
methodus directa, quae a solutione aequationis $axx+2bxy+cyy=M$
omnino independens esset.\ Hanc iam ipsi sequenti articulo reservamus.

\subsubsection*{330.}

Quamquam haec methodus indirecta, proposita in art.\ 327, variis
modis perfici atque contrahi possit, tamen in casibus quibusdam difficilior
manere potest, quam ut usui practico commode inserviat.\ Notandum autem
est, hanc methodum potissimum eo spectare, ut valores expr.\ $\sqrt{-D}$ (mod.\ $M$)
inveniantur pro modulis $M$ maximis, methodoque directa art.\ 322 vix
adhiberi possint; pro modulis minoribus methodus directa semper est
praeferenda.

\subsection*{Duae methodi numerorum factors investigandi.}

\subsubsection*{331.}

Iam ad aliam gravissimam quaestionem progredimur, scilicet ad
factoring numerorum.\ Quamquam vero in arithmetica elementarii plures
methodi traduntur, quibus numeri compositi in factores suos primos resolvi
possint, attamen hae omnes tam lentae sunt, ut vix ad numeros aliquot
notarum aptae esse videantur.\ Methodi autem, quae in disquisitionibus
praecc.\ expositae sunt, variis modis ad hoc problema applicari possunt,
quorum primum hic exponemus, alterumque in articulo sequente.

Methodus prior eo nititur, quod si numerus propositus $M$ ex pluribus
numeris primis (incl.\ eorundem potestatibus) compositus est, maiorem quam
$\sqrt{M}$ factorem habere nequit, nisi unus tantum aliorum factorum primorum
esset.\ Quodsi itaque omnes numeri primi infra $M$ (exclusis iis, per quos
divisio frustra iam tentata est) per $Q$ designentur, manifesto omnes divisores
primi ipsius $M$, si quos habet, in $Q$ contenti erunt.\ Iam si aliquis numerus
$r$ (non-quadratum) esse residuum quadraticum ipsius $M$, omnes numeri
primi, quorum $r$ est non-residuum, certo numerum $M$ metiri non poterunt;
quare ex $Q$ omnes huiusmodi numeros primos (qui plerumque semissem fere
efficient) eiicere licebit.

Si insuper de alio numero non-quadrato $r'$ constat, ipsum esse residuum
ipsius $M$, e numeris primis in $Q$ post primam exclusionem relictis iterum
eos excludere poterimus, quorum NR.\ est $r'$, qui rursus illorum semissem fere
conficient, siquidem residua $r$ et $r'$ sunt independentia, i.\ e.\ nisi alterum
necessario per se est residuum omnium numerorum, quorum residuum est
alterum, quod eveniret, quando $rr'$ esset quadratum.

Si adhuc alia residua ipsius $M$ noti sunt, $r''$, $r'''$ etc., quae omnia a
reliquis sunt independentia, cum singulis exclusiones similes institui possunt,
per quas multitudo numerorum in $Q$ rapidissime diminuetur, ita ut mox vel
omnes deleti sint, in quo casu $M$ certo erit numerus primus; vel tam pauci
restant (inter quos omnes divisores primi ipsius $M$, si quos habet, manifesto
reperientur), ut divisio per ipsos nullo negotio tentari possit.

Pro numero millionem non multum superante plerumque sex aut septem
exclusiones abunde sufficient; pro numero ex octo aut novem figuris
constante novem aut decem exclusiones abunde sufficient.\ Duo iam sunt,
de quibus agere oportebit, primo quomodo residua ipsius $M$ idonea et satis
multa inveniri possint, deinde quo pacto exclusionem ipsam commodissime
perficere liceat.

\subsubsection*{332.}

Numeros primos, quorum residuum est numerus datus $r$ (quem per
nullum quadratum divisibilem supponere liceat), ab iis, quorum non-residuum
est, sive divisores expr.\ $xx-r$ a non-divisoribus distinguere, in Sectione IV
copiose docuimus, scilicet omnes priores sub certis huiusmodi formulis
$4rz+a$, $4rz+b$ etc.\ aut talibus $rz+a$, $rz+b$ etc.\ contentos esse,
posterioresque sub aliis similibus.

Quoties $r$ est numerus valde parvus, exclusiones adiumento harum
formularum percommode perfici possunt; e.\ g.\ excludendi erunt omnes
numeri formae $4z+1$, $4z+2$, $4z+3$, quando $r=-1$; omnes numeri
formarum $8z+3$ et $8z+5$, quando $r=2$; omnes numeri formarum
$8z+5$ et $8z+7$, quando $r=-2$ etc.

Sed quum non semper in potestate sit, huiusmodi residua numeri
propositi invenire, neque formularum applicatio pro valore magno ipsius $r$
satis commoda sit, ingens lucrum est, laboremque exclusionis mirifice
sublevat, si pro multitudine satis magna numerorum per quadratum non
divisibilium tum positivorum tum negativorum tabula iam constructa habetur,
in qua numeri primi, quorum residua sunt illi singuli, ab iis, quorum
non-residua sunt, distinguuntur.\ Talis tabula perinde adornari poterit ac
specimen ad calcem huius operis adiectum supraque iam descriptum; sed ut
ad institutum praesens utilitatem satis amplam praestet, numeri primi in
margine positi (moduli) longe ulterius puta saltem usque ad 1000 aut ad
10000 continuati esse debent, praetereaque commoditas multum augetur, si
in facie etiam numeri compositi et negativi recipiuntur, etsi hoc non sit
absolute necessarium, ut e Sectione IV perspicuum est.\ Ad summum autem
commoditatis fastigium usus talis tabulae evehetur, si singulae columellae
verticales, e quibus constat, exsecantur lamellisque aut baculis (Neperianis
similibus) agglutinantur, ita ut eae, quae in quovis casu sunt necessariae
i.\ e.\ quae numeris $r$, $r'$, $r''$ etc., residuis, respondent, separate examinari
possint.

Quibus iuxta tabulae columnam primam (quae modulos exhibet) rite
positis i.\ e.\ ita ut loca singulorum baculorum eidem numero primo columnae
primae respondentium cum hoc in directum iaceant, sive in eadem linea
horizontali siti sint: manifesto ei numeri primi, qui post exclusiones cum
residuis $r$, $r'$, $r''$ etc.\ ex $Q$ remanent, per solam inspectionem immediate
cognosci poterunt, nimirum hi convenient cum iis in columna prima quibus
in omnibus baculis adiacentibus lineolae respondent, reiicique debent omnes,
quibus in ullo bacillo spatium vacuum adiacet.

Si alicunde constat, numeros $-6$, $+13$, $-14$, $+17$, $+37$, $-53$
esse residua ipsius 997331, consociandae erunt columna prima (quae in hoc
casu usque ad 997 continuata esse debet, i.\ e.\ usque ad numerum primum
proxime minorem quam $\sqrt{997331}$) atque lamellae, in quarum facie numeri
$-6$, $+13$ etc.\ sunt suprascripti.

\subsubsection*{333.}

Modus residua ipsius $M$ eruendi duplici potissimum via institui potest.

\textbf{I.}\ Methodus simplicissima, iisque, qui per frequentem exercitationem
iam aliquam dexteritatem sibi conciliaverunt, commodissima, consistit in eo,
ut $M$ quomodocunque in duas partes decomponatur $kM=a+b$ (sive utraque
sit positiva sive altera positiva altera negativa), quarum productum signo
mutato erit residuum ipsius $M$; erit enim $-ab\equiv aa\equiv bb$ (mod.\ $M$)
adeoque $-ab$ R$M$.\ Numeri $a$, $b$ ita accipiendi sunt, ut productum per
quadratum magnum divisibile quotiensque vel parvus vel saltem in factores
non nimis magnos resolubilis evadat, quod semper non difficile effici poterit.
Imprimis commendandum est, ut pro $a$ accipiatur vel quadratum, vel
quadratum duplex, vel triplex etc.\ a numero $M$ numero vel parvo vel in
factores commodos resolubili discrepans.\ Ita e.\ g.\ invenitur
\begin{align*}
997331 &= 999^2-2\cdot 5\cdot 67 = 994^2+5\cdot 11\cdot 13^2 \\
&= 2\cdot 706^2+3\cdot 17\cdot 3^2 = 3\cdot 575^2+11\cdot 31\cdot 4^2 \\
&= 3\cdot 577^2-7\cdot 13\cdot 4^2 = 3\cdot 578^2-7\cdot 19\cdot 37 \\
&= 11\cdot 299^2+2\cdot 3\cdot 5\cdot 29\cdot 4^2 = 11\cdot 301^2+5\cdot 12^2.
\end{align*}
Hinc habentur residua sequentia $2\cdot 5\cdot 67$, $-5\cdot 11$, $-2\cdot 3\cdot 17$,
$-3\cdot 11\cdot 31$, $3\cdot 7\cdot 13$, $3\cdot 7\cdot 19\cdot 37$, $-2\cdot 3\cdot 5\cdot 11\cdot 29$; discerptio
ultima suppeditat residuum $-5\cdot 11$ quod iam habemus.\ Pro residuis
$-3\cdot 11\cdot 31$, $-2\cdot 3\cdot 5\cdot 11\cdot 29$ haec adoptare possumus $3\cdot 5\cdot 31$,
$2\cdot 3\cdot 29$, ex illorum combinatione cum $-5\cdot 11$ oriunda.

\textbf{II.}\ Methodus secunda et tertia inde petuntur, quod, si duae formae
binariae $(A,B,C)$, $(A',B',C')$ eiusdem determinantis $D$, aut $-D$, aut
generalius $\pm kM$, ad idem genus pertinent, numeri $AA'$, $AC$, $A'C$,
$A'C'$ sunt residua ipsius $kM$; hoc nullo negotio inde perspicitur, quod
numerus quivis characteristicus unius formae, puta $m$, etiam est numerus
char.\ alterius, adeoque $mA$, $mC$, $mA'$, $mC'$ omnes residua ipsius $kM$.
Si itaque $(a,b,a')$ est forma reducta determinantis $M$ aut generalius
$kM$, atque $(a',b',a'')$, $(a'',b'',a''')$ etc.\ formae ex ipsius periodo tentae:
quae, adeoque ipsi aequivalentes et a potiori sub eodem genere continentur,
numeri $aa'$, $aa''$, $aa'''$ etc.\ omnes erunt residua ipsius $M$.

Computus multitudinis magnae formarum talis periodi facillime adiumento
algorithmi art.\ 187 instituitur; quae factores nimis magnos implicant, erunt
reiicienda.\ Ea quae factores simplicissima plerumque prodeunt statuendo
$a=1$.\ Ecce initia periodorum formarum $(1,998,1327)$ et
$(1,1412,-918)$, quarum determinantes sunt $997331$, $1994662$:

\begin{center}
\begin{tabular}{ll}
$(1,998,1327)$ & $(1,1412,-918)$ \\
$2\mathfrak{C}\ (3,997,-1326)$ & $2\mathfrak{C}\ (2,1412,-1377)$ \\
$3\mathfrak{C}\ (27,976,-7)$ & $3\mathfrak{C}\ (9,1386,-484)$ \\
$4\mathfrak{C}\ (81,947,-14)$ & $4\mathfrak{C}\ (27,1342,-163)$ \\
$5\mathfrak{C}\ (243,892,-21)$ & $5\mathfrak{C}\ (81,1236,-121)$ \\
$6\mathfrak{C}\ (729,781,-28)$ & $6\mathfrak{C}\ (243,1027,-142)$ \\
$7\mathfrak{C}\ (2187,476,-63)$ & $7\mathfrak{C}\ (729,342,1085)$ \\
$8\mathfrak{C}\ (1027,209,\quad\,\,\,2187)$ & $8\mathfrak{C}\ (1027,342,1085)$ \\
$9\mathfrak{C}\ (932,-437,1275)$ & $9\mathfrak{C}\ (1085,\,\,\,\dots)$
\end{tabular}
\end{center}

Hinc proraanant residua (inutilibus reiectis) $3\cdot 476$, $1027$, $1085$,
$425$ sive (tolendo factores quadratos) $3\cdot 7\cdot 17$, $13\cdot 79$, $5\cdot 7\cdot 31$,
$17$, e quorum combinatione apta cum octo residuis in II inventis facile
eruuntur duodecim sequentia:

$2\cdot 3$, $13$, $-2\cdot 7$, $-5\cdot 11$, $79$, $-83$, $-29$, $-2\cdot 59$,
$-2\cdot 5\cdot 31$, $2\cdot 5\cdot 67$, $17$, $37$, $-53$;

sex priora sunt eadem, quibus in art.\ 331 usi sumus.

Adiici potuissent residua 19 et 53, quae in I reperta sunt; reliqua
eruta ab iis quae hie evolvimus iam sunt dependentia.

\section*{De aequationibus circuli sectiones definientibus.}

\subsection*{Aequationes pro functionibus trigonometricis arcuum, qui sunt
pars aut partes totius peripheriae; reductio functionum trigonometricarum
ad radices aequationis $x^n-1=0$.}

\subsubsection*{337.}

Satis constat, functiones trigonometricas omnium angulorum
$\frac{kP}{n}$, denotando per $k$ indefinite omnes numeros $0$, $1$, $2\dots n-1$,
per radices aequationum $n^{ti}$ gradus exprimi, puta sinus per radices huius (I)
\begin{multline*}
x^n - \frac{n}{1}x^{n-1} + \frac{n(n-3)}{1\cdot 2}x^{n-2} - \frac{n(n-4)(n-5)}{1\cdot 2\cdot 3}x^{n-3} \\
+ \text{etc.}\ \pm \frac{n}{1}x = 0
\end{multline*}

Cosinus per radices huius (II)
\begin{multline*}
x^n - nx^{n-2} + \frac{n(n-3)}{1\cdot 2}x^{n-4} - \frac{n(n-4)(n-5)}{1\cdot 2\cdot 3}x^{n-6} \\
+ \text{etc.}\ \pm 2x - 2\cos\frac{2kP}{n} = 0
\end{multline*}

Hae aequationes (quae generaliter pro quovis valore impari ipsius $n$
valent; scilicet I et III, dividendo partem a laeva per $x$; II vero pro pari
quoque), ponendo $n=2m+1$, facile ad gradum $m^{ti}$ deprimi possunt.
Aequatio II autem manifesto radicem $x=2$ ($=2\cos 0$) implicat, et e
reliquis aequationis $x^n-1=0$ partibus, dividendo per $x-1$ et
substituendo $y$ pro $xx$, quadratam extrahendo, aequatio II reducitur ad
hanc
\[y^m - \frac{n}{2}y^{m-1} + \frac{n(n-2)}{16}y^{m-2} - \frac{n(n-3)(n-4)}{16\cdot 32}y^{m-3} + \text{etc.}\]
cuius radices erunt cosinus angulorum $\frac{P}{n}$, $\frac{2P}{n}$, $\frac{3P}{n}\dots\frac{mP}{n}$.

Ulteriores reductiones harum aequationum, pro eo quidem casu, ubi $n$
est numerus primus, hactenus non habebantur.

Attamen nulla harum aequationum tam tractabilis et ad institutum
nostrum tam idonea est, quam haec $x^n-1=0$, cuius radices cum radicibus
illarum arctissime connexas esse constat.\ Scilicet, scribendo brevitatis
caussa $i$ pro quantitate imaginaria $\sqrt{-1}$, radices aequationis $x^n-1=0$
exhibentur per
\[\cos\frac{kP}{n} + i\sin\frac{kP}{n}\]
ubi pro $k$ accipiendi sunt omni numeri $0$, $1$, $2\dots n-1$.\ Quocirca
quum sit
\[\cos\frac{kP}{n} + i\sin\frac{kP}{n} = \left(\cos\frac{P}{n} + i\sin\frac{P}{n}\right)^k\]
radices aequationis I exhibebuntur per
\[\frac{1}{2i}\left(\left(\cos\frac{P}{n} + i\sin\frac{P}{n}\right)^k - \left(\cos\frac{P}{n} - i\sin\frac{P}{n}\right)^k\right)\]
sive per
\[\frac{x^k - x^{-k}}{2i}\]
radices aequationis II per
\[\frac{x^k + x^{-k}}{2}\]
denique radices aequationis III per $\frac{x^k - x^{-k}}{x^k + x^{-k}}$, sive per
\[\frac{x^{2k} - 1}{x^{2k} + 1}\]
Hanc ob caussam disquisitionem considerationi aequationis $x^n-1=0$
superstruemus, ipsum $n$ esse numerum primum imparem supponendo.\ Ne
vero investigationum ordinem interrumpere oporteat, sequens lemma hie
praemittimus.

\subsubsection*{338.}

\textbf{Problema.}\ Data aequatione
\[x^m + Ax^{m-1} + \text{etc.} = 0\quad (W)\]
invenire aequationem ($W'$), cuius radices sint potestates $m^{tae}$ radicum
aequationis ($W$), designante $\lambda$ exponentem integrum positivum datum.

\textbf{Solutio.}\ Designatis radicibus aequationis $W$ per $a$, $b$, $c$ etc.,
radices aequ.\ $W'$ invenire licet per $a^\lambda$, $b^\lambda$, $c^\lambda$ etc.\ Per theorema
notum \textsc{Newtonianum} e coefficientibus aequ.\ $W$ aggregata quarumlibet
potestatum radicum $a$, $b$, $c$ etc.\ deduci possunt.\ Quaerantur itaque
summae $a^\lambda+b^\lambda+c^\lambda+$ etc., $a^{2\lambda}+b^{2\lambda}+c^{2\lambda}$ etc.\ etc.
umdem theorema coefficientes aequ.\ $W'$ via inversa per $a^\lambda+b^\lambda+c^\lambda+$
etc., $a^{2\lambda}+b^{2\lambda}+c^{2\lambda}$ etc.\ deduci poterunt.\ Q.\ E.\ F.

Simul hinc liquet, si omnes coefficientes in $W$ sunt rationales, omnes
quoque in $W'$ rationales evadere.\ Alia quidem via probari potest, si illi
omnes integri sint, etiam hos omnes integros fieri; huic autem theoremati,
ad institutum nostrum non adeo necessario, hie non immoramur.

\subsubsection*{339.}

Aequatio $x^n-1=0$ (in suppositione, quae semper abhinc
subintelligenda est, $n$ esse numerum primum imparem) unicam radicem
realem implicat $x=1$ sive $x-1=0$; reliquae, quas aequatio
$x^{n-1}+x^{n-2}+\text{etc.}+x+1=0$ complectitur, omnes sunt imaginariae;
harum complexum per $Q$, functionemque $x^{n-1}+x^{n-2}+\text{etc.}+x+1$
per $X$ denotabimus.

Si itaque $r$ est radix quaecunque ex $Q$, erit $r^n=1$, et generaliter
$r^{\lambda n}=1$ pro quovis valore integro ipsius $\lambda$, hinc perspicuum est, si
$\lambda$, $\mu$ sint integri secundum mod.\ $n$ incongrui, fore $r^\lambda$, $r^\mu$ inaequales;
vero si $\lambda$, $\mu$ secundum mod.\ $n$ congrui, ita accipi potest, ut
$\lambda-\mu = vn$, adeoque $r^{\lambda-\mu} = r^{vn} = 1$, casu integer $v$ positive seu
negativo; unde $r^\lambda = r^\mu$, i.\ e.\ $r^\lambda$, $r^\mu$ convenire.\ Cum $r^{n-1} = \frac{1}{r}$,
$1$, $r$, $rr$, $r^3\dots r^{n-1}$ convenire cum $1$, $r$, $rr$, $r^3\dots r^{n-1}$.

Si omnes $1$, $r$, $rr$, $r^3\dots r^{n-1}$ sunt diversae, hae exhibebunt omnes
radices aequ.\ $x^n-1=0$.\ Si vero $r^e=1$, i.\ e.\ $r$ est radix aequ.
$x^e-1=0$, ubi $e$ integer quicunque per $n$ non divisibilis, Erit itaque
$r^e=1$, et proin $1+r^e+r^{2e}+r^{3e}\dots+r^{(n-1)e}=n$, unde
$1+r^e+r^{2e}+r^{3e}\dots+r^{(n-1)e} = n$ vel $=0$, prout $\lambda$ per $n$
divisibilis est vel divisibilis.

Duas radices tales ut $r$ et $\frac{1}{r}$ ($=r^{n-1}$), aut generaliter $r^\lambda$ et
$r^{-\lambda}$ \emph{reciprocas} vocabimus; manifestum est, productum ex duobus
factoribus simplicibus $x-r$ et $x-\frac{1}{r}$ fieri $=xx-2x\cos\omega+1$, ita ut
angulus $\omega$ vel angulo $\frac{\lambda P}{n}$ vel alicui multiplo eius sit aequalis.

\subsubsection*{340.}

Quoniam itaque, una radice ex $Q$ per $r$ expressa, omnes radices
aequ.\ $x^n-1=0$ per potestates ipsius $r$ exprimuntur, productum, e
pluribus radicibus huius aequ.\ quomodocunque conflatum, per $r^\lambda$
exhiberi poterit, ita ut $\lambda$ sit vel 0, vel positivus et $<n$.\ Designando
itaque per $\varphi(t,u,v\dots)$ functionem algebraicam rationalem integram
indeterminatarum $t$, $u$, $v$ etc.\ talem, ut si pro $t$, $u$, $v$ etc.\ quaedam e
radicibus aequ.\ $x^n-1=0$ substituantur, puta $t=a$, $u=b$, $v=c$
etc., haec sub formam $A+A'r+A''rr+A'''r^3\dots+A^{(n-1)}r^{n-1}$
reduci posse, ita ut coefficientes $A$, $A'$ etc.\ (e quibus etiam aliqui deesse
adeoque $=0$ fieri possunt) sint quantitates determinatae, insuperque omnes
hos coefficientes integros fieri, si omnes coefficientes determinati in
$\varphi(t,u,v\dots)$, i.\ e.\ omnes hi $h$ sint integri.

Quodsi vero postea pro $t$, $u$, $v$ etc.\ substituuntur $aa$, $bb$, $cc$
etc., quae antea reducebantur ad $r^\lambda$, nunc fiet $r^{2\lambda}$, unde facile concluditur,
fieri
\[\varphi(aa,bb,cc\dots) = A+A'rr+A''r^4\dots+A^{(n-1)}r^{2(n-1)}.\]
Perinde erit generaliter, pro valore quocunque integro ipsius $\lambda$,
\[\varphi(a^\lambda,b^\lambda,c^\lambda\dots) = A+A'r^\lambda+A''r^{2\lambda}\dots+A^{(n-1)}r^{(n-1)\lambda}\]
quae propositio maximi est momenti fundamentumque disquisitionum
sequentium constituit.

Hinc sequitur etiam
\[\varphi(1,1,1\dots) = nA = A+A'+A''\dots+A^{(n-1)}\]
nee non
\begin{multline*}
\varphi(a,b,c\dots) + \varphi(aa,bb,cc\dots) + \varphi(a^3,b^3,c^3\dots) + \dots \\
+ \varphi(a^{n-1},b^{n-1},c^{n-1}\dots) = nA
\end{multline*}
quae itaque summa semper fit integra per $n$ divisibilis, quando omnes
coefficientes determinati in $\varphi(t,u,v\dots)$ sunt integri.

\subsection*{Theoria radicum aequationis $x^n-1=0$ (ubi supponitur $n$
esse numerum primum).}

\subsubsection*{341.}

\textbf{Theorema.}\ Si functio $X$ per functionem inferioris gradus
\[P = x^\lambda + Ax^{\lambda-1} + Bx^{\lambda-2}\dots + Kx + L\]
est divisibilis, coefficientes $A$, $B\dots L$ omnes integri esse nequeunt.

\textbf{Dem.}\ Sit $X=PQ$, atque $\mathfrak{P}$ complexus radicum aequationis
$P=0$, $\mathfrak{Q}$ complexus radicum aequationis $Q=0$, ita ut $Q$ constet
ex $\mathfrak{P}$ et $\mathfrak{Q}$ simul sumtis.\ Porro sit $\mathfrak{R}$ complexus radicum
ipsis $\mathfrak{P}$ reciprocarum, $\mathfrak{S}$ complexus radicum ipsis $\mathfrak{Q}$
reciprocarum, sintque radices, quae continentur in $\mathfrak{R}$, eaeque quae
continentur in $\mathfrak{S}$, radices aequationis $R=0$, $S=0$, facile
perspicitur, si $\mathfrak{R}$ et $\mathfrak{S}$ radices aequationis $R=0$, $S=0$, facile
perspicitur, si $\mathfrak{R}$ et $\mathfrak{S}$ iunctae complexum $Q$ efficient, ac erit
$RS=X$.

Iam quatuor casus distinguimus.

\textbf{I.}\ Quando $\mathfrak{P}$ convenit cum $\mathfrak{R}$ adeoque $P=R$.\ In hoc
casu manifesto binae semper radices in $\mathfrak{P}$ reciprocae erunt, adeoque $P$
productum ex $\frac{1}{2}\lambda$ factoribus talibus duplicibus $xx-2x\cos\omega+1$; quum
talis factor sit $=(x-\cos\omega)^2+\sin\omega^2$, $P$ pro valore quocunque reali
ipsius $x$ necessario valorem realem positivum obtinere.

Sint aequationes, quarum radices sunt quadrata, cubi, biquadrata
$\dots$ $n^{tae}$ potestates radicum in $\mathfrak{P}$ resp.\ hae $P'=0$, $P''=0$,
$P'''=0$ etc., sintque valores functionum $P$, $P'$, $P''$ etc., quos
obtinent statuendo $t$, $u$, $v$ etc.\ radices in $\mathfrak{P}$, resp.\ $p$, $p'$,
$p''$ etc., itaque $p$ sit valor functionis $P$; erit per art.\ 340
$p+p'+p''\dots+p^{(n-1)} = nA$, adeoque summa $p+p'+p''\dots+p^{(n-1)}$
integer per $n$ divisibilis.\ Praeterea facile perspicietur, si $P=(1-t)(1-u)(1-v)$
etc., valor eiusdem, statuendo pro $t$, $u$, $v$ etc.\ quadrata illarum
radicum etc.\ manifesto $P'=(1-tt)(1-uu)(1-vv)$ etc., insuperque
valor pro $t=u=v=1$, $P=0$, tunc per ante dicta $p$ erit quantitas
positiva et prorsus simili ratione etiam $p'$, $p''$ etc.\ positivae erunt.

\textbf{II.}\ Quando $\mathfrak{P}$ et $\mathfrak{Q}$ nullam radicem communem habent,
omnes radices $\mathfrak{R}$ in $\mathfrak{P}$, omnesque $\mathfrak{S}$ in $\mathfrak{Q}$ necessario
reperientur, unde erit $P=R$ et $Q=S$.\ Quamobrem in $X=PQ$ erit
productum ex $P$ et $Q$, sive $x^n-1$ divisibilis per $xx-(A-L^{-1})x+1$,
unde statuendo $x=1$, fit $n$ divisibilis per $2-A-L^{-1}$, adeoque
$L^{-1}$ qui coefficientem ultimum in $X$ metiri deberet, necessario foret
$=\pm 1$, unde $\pm n$ esset numerus quadratus.\ Quod quum hypothesi
repugnet, suppositio consistere nequit.

\textbf{III.}\ Quando $\mathfrak{P}$ et $\mathfrak{Q}$ vel coincidunt, vel saltem radices
communes implicant, prorsus eodem modo omnes coefficientes in $Q$
rationales esse nequeunt; fierent vero rationales, si omnes in $P$ adeoque
etiam omnes in $P$ rationales essent; hoc itaque est impossibile.

\textbf{IV.}\ Si vero neque $\mathfrak{P}$ cum $\mathfrak{R}$, neque $\mathfrak{Q}$ cum $\mathfrak{S}$
ullam radicem communem habet, omnes radices $\mathfrak{P}$ in $\mathfrak{S}$, omnesque
$\mathfrak{Q}$ in $\mathfrak{R}$ necessario reperientur, unde erit $P=S$ et $Q=R$.

Quamobrem in $X=PQ$ erit productum ex $x^\lambda+Ax^{\lambda-1}+\dots+L$
et $x^\mu+A'x^{\mu-1}+\dots+L'$, unde statuendo $x=1$, fit $n=LL'$.
Iam si omnes coefficientes in $P$ rationales, adeoque per art.\ 42 etiam
integri essent, $L$ qui coefficientem ultimum in $X$ metiri deberet,
necessario foret $=\pm 1$, unde $L'$ esset $=\pm n$, et proin $Q$ radicem
$x=\pm n$ haberet, quod fieri nequit.\ Quod quum hypothesi repugnet,
suppositio consistere nequit.

Ex hoc itaque theoremate liquet, quomodocunque $X$ in factores
resolvatur, horum coefficientes partim saltem irrationales fieri, adeoque
aliter, quam per aequationem elevatam, determinari non posse.

\subsubsection*{342.}

\textbf{Propositum disquisitionum sequentium,} quod paucis declaravisse
haud inutile erit, eo tendit, ut $X$ in factores continuo plures gradatim
resolvatur, et quidem ita, ut horum coefficientes per aequationes ordinis
quam infimi determinentur, usque dum hoc modo ad factores simplices
sive ad radices $Q$ ipsas perveniatur.

Scilicet ostendemus, si numerus $n-1$ quomodocunque in factores
integros $\alpha$, $\beta$, $\gamma$ etc.\ resolvatur (pro quibus singulis numeros primos
accipere licet), $X$ in $\alpha$ factores $\beta\gamma$ dimensionum resolvi posse, quorum
coefficientes per aequationem $\alpha^{ti}$ gradus determinentur; singulos hos
factores iterum in $\beta$ alios $\gamma$ dimensionum adiumento aequationis
$\beta^{ti}$ gradus etc., ita ut designante $\nu$ multitudinem factorum
$\alpha$, $\beta$, $\gamma$ etc., inventio radicum $Q$ ad resolutionem $\nu$ aequationum
$\alpha^{ti}$, $\beta^{ti}$, $\gamma^{ti}$ etc.\ gradus reducatur.

E.\ g.\ pro $n=17$, ubi $n-1=2\cdot 2\cdot 2\cdot 2$, quatuor aequationes
quadraticas solvere oportebit; pro $n=73$, ubi $n-1=2\cdot 2\cdot 2\cdot 3\cdot 3$,
tres quadraticas duasque cubicas.

Quum in sequentibus persaepe tales potestates radicis $r$ considerandae
sint, quarum exponentes rursus sunt dignitates, huiusmodi expressiones autem
non sine molestia typis describantur: ad facilitandam impressionem sequenti
in posterum abbreviatione utemur.

Pro $r$, $rr$, $r^3$ etc.\ scribemus $[1]$, $[2]$, $[3]$ etc., generaliterque pro
$r^\lambda$, denotante $\lambda$ integrum quemcunque, $[\lambda]$.\ Tales itaque expressiones
penitus determinatae nondum sunt, sed fiunt, simulac pro $r$ sive $[1]$ radix
determinata ex $Q$ accipitur.\ Erunt itaque generaliter $[\lambda]$, $[\mu]$ aequales
vel inaequales, prout $\lambda$, $\mu$ secundum modulum $n$ congrui sunt vel
incongrui; porro $[0]=1$; $[\lambda]\cdot[\mu]=[\lambda+\mu]$; $[\lambda^\nu]=[\lambda\nu]$; summa
$[0]+[\lambda]+[2\lambda]\dots+[(n-1)\lambda] = n$ vel $=0$, prout $\lambda$ per $n$
divisibilis est vel divisibilis.

\subsection*{Omnes radices $Q$ in certas classes (periodos) distribuuntur.}

\subsubsection*{343.}

Si, pro modulo $n$, $g$ est numerus talis, qualem in Sectione III
radicem primitivam diximus, $n-1$ numeri $1$, $g$, $gg\dots g^{n-2}$ secundum
mod.\ $n$ congrui erunt, ordine tantum excepto, numeris $1$, $2$, $3\dots n-1$.
Hinc sponte sequitur, radices $[1]$, $[g]$, $[gg]\dots[g^{n-2}]$ cum $Q$
coincidere et prorsus simili modo generalius $[\lambda]$, $[\lambda g]$, $[\lambda gg]\dots[\lambda g^{n-2}]$
cum $Q$ coincidere, designante $\lambda$ integrum quemcunque per $n$ non
divisibilem.

Si itaque $G$ est alia radix primitiva, radices $[1]$, $[G]$, $[GG]\dots[G^{n-2}]$
etiam cum $Q$ convenient, facile probatur, si $G=g^\mu$, ubi $\mu$ est divisor
ipsius $n-1$, atque ponatur $n-1=ef$, $g^e=H$, etiam $f$ numeros
$H$, $H^2$, $H^3\dots H^{f-1}$ secundum mod.\ $n$ congruos esse (sine respectu
ordinis).\ Supponamus enim $G^{\mu\nu}\equiv H^\nu$ (mod.\ $n$), sitque $\mu$ numerus
arbitrarius positivus et $<f$ atque $\nu$ residuum minimum ipsius
$\mu\omega$ (mod.\ $f$).\ Tunc erit $G^{\mu\nu}=g^{\mu\nu\omega}=g^{\mu\omega ef}=g^{\nu ef}\equiv H^\nu$
(mod.\ $n$), i.\ e.\ $g^{\mu\nu}\equiv H^\nu$, hinc quivis numerus posterioris seriei
alicui in serie $1$, $H$, $HH\dots H^{f-1}$ congruum habebit, et perinde vice
versa.\ Hinc manifestum est, radices $[\lambda]$, $[\lambda H]$, $[\lambda HH]\dots[\lambda H^{f-1}]$
identicas esse cum his $[\lambda]$, $[\lambda g^e]$, $[\lambda g^{2e}]\dots[\lambda g^{(f-1)e}]$, generaliusque
eodem modo facile perspicietur,
\[[\lambda], [\lambda h], [\lambda hh]\dots[\lambda h^{f-1}]\]
cum
\[[\lambda], [\lambda g^e], [\lambda g^{2e}]\dots[\lambda g^{(f-1)e}]\]
convenire.\ Aggregatum talium $f$ radicum
\[[\lambda]+[\lambda h]+[\lambda hh]+\text{etc.}+[\lambda h^{f-1}]\]
quod, quum non mutetur accipiendo pro $g$ aliam radicem primitivam
quam $\mathfrak{g}$ considerandum est, per $(f,\lambda)$ vocabimus, tamquam independens
earundem radicum complexum designabimus; ubi ad radicum ordinem non
respicitur.\ *)

In exhibenda tali periodo e re erit, singulas radices, e quibus constat,
ad expressionem simplicissimam reducere, puta pro numeris
$\lambda$, $\lambda h$, $\lambda hh$ etc.\ residua minima sec.\ mod.\ $n$ substituere, secundum
quorum magnitudinem, si placet, etiam partes periodi ordinari poterunt.

\textbf{E.\ g.}\ Pro $n=19$, ubi 2 est radix primitiva, periodus $(6,1)$ constat
e radicibus $[1]$, $[8]$, $[64]$, $[512]$, $[4096]$, $[32768]$, sive $[1]$, $[7]$,
$[8]$, $[11]$, $[12]$, $[18]$.\ Similiter periodus $(6,3)$ constat ex $[3]$,
$[14]$, $[2]$, $[5]$, $[16]$, $[17]$.\ Periodus $(6,2)$ cum praec.\ identica
invenitur.\ Periodus $(6,4)$ continet $[4]$, $[6]$, $[9]$, $[10]$, $[13]$, $[15]$.

\textit{Varia theoremata de periodis radicum $Q$.}

\subsubsection*{344.}

Circa huiusmodi periodos statim se offerunt observationes sequentes:

\textbf{I.}\ Quum sit $\lambda h^f\equiv\lambda$, $\lambda h^{f+1}\equiv\lambda h$ etc.\ (mod.\ $n$),
manifestum est, ex iisdem radicibus, e quibus constet $(f,\lambda)$, etiam
constare $(f,\lambda')$, generaliter itaque designante $[\lambda']$ radicem quamcunque ex
$(f,\lambda)$, haec periodus $(f,\lambda')$ cum $(f,\lambda)$ omnino identica erit.\ Si itaque
duae periodi ex aeque multis radicibus constantes (quales \emph{similes} dicemus)
ullam radicem communem habent, manifesto identicae erunt.

Quare fieri nequit, ut duae radices in aliqua periodo simul contineantur,
in alia simili vero una earum tantum reperiatur; porro patet, si duae
radices $[\lambda]$, $[\lambda']$ in eadem periodum $f$ terminorum pertineant, valorem
expr.\ $\frac{\lambda'}{\lambda}$ (mod.\ $n$) ipsius potestati $h^e$ congruum esse, sive supponi
posse $\lambda'\equiv\lambda h^e$ (mod.\ $n$), $e=0$, $1$, $2\dots f-1$.

\textbf{II.}\ Si $n-1$ est productum e tribus integris positivis $a$, $b$, $c$,
manifestum est, quamvis periodum $abc$ terminorum compositam esse, e
periodis $c$ terminorum contentarum in $(bc,\lambda)$, $(bc,\lambda g^a)$,
$(bc,\lambda g^{2a})\dots(bc,\lambda g^{(b-1)a})$, unde hae sub $(ac,\lambda)$, $(ac,\lambda g^b)$,
$(ac,\lambda g^{2b})\dots(ac,\lambda g^{(a-1)b})$ contentae dicentur.

Ita e.\ g.\ pro $n=19$, $a=3$, $b=3$, $c=2$, $Q$ constat e tribus
periodis $(6,1)$, $(6,2)$, $(6,4)$, ad quarum aliquam quaevis alia similis
reducitur.\ Ita $(6,0)$ manifesto ex unitatibus est composita.

Aeque facile perspicitur, si $n-1$ est productum e quatuor integris
$a$, $b$, $c$, $d$, periodus $abcd$ terminorum composita erit e $b$ periodis
$ac$ terminorum ex $c$ periodis $d$ terminorum etc.\ In universum si
$n-1=\alpha\beta\gamma\dots$, periodus $n-1$ terminorum, i.\ e.\ $Q$ ipsa,
composita erit e $\alpha$ periodis $\beta\gamma\dots$ terminorum, quaeque harum iterum
ex $\beta$ periodis $\gamma\dots$ terminorum etc.

\textbf{III.}\ Quum quaevis periodus, cuius terminorum multitudo est par
$=2a$, in $a$ periodos binorum terminorum discerpi possit, patet generalius,
valorem cuiusvis aggregati, in quo terminorum multitudo par est, esse
quantitatem realem.\ Hinc facile deducitur, valorem producti e duabus
periodis, in quibus terminorum multitudo est par, semper quoque realem
esse.

\subsubsection*{345.}

\textbf{Theorema.}\ Sint $(f,\lambda)$, $(f,\mu)$ duae periodi similes, identicae aut
diversae, constetque $(f,\lambda)$ e radicibus $[\lambda]$, $[\lambda']$, $[\lambda'']$ etc., erit
aggregatum $(f,\lambda)\cdot(f,\mu)$ productum ex $f$ periodorum similium, puta
\[W = (f,\lambda+\mu) + (f,\lambda'+\mu) + (f,\lambda''+\mu) + \text{etc.}\]

\textbf{Dem.}\ Sit ut supra $n-1=cf$; $g$ radix primitiva pro modulo $n$,
atque $g^e=h$, unde per praecedentia erit
\[(f,\lambda) = (f,\lambda h) = (f,\lambda hh)\text{ etc.}\]
Hinc productum quaesitum erit
\[[\lambda]+[\lambda']+[\lambda'']\text{ etc.}\cdot[\mu]+[\mu h]+[\mu hh]\text{ etc.}\]
adeoque
\[[\lambda+\mu]+[\lambda'+\mu]+[\lambda''+\mu]+\text{etc.}\]
\[+ [\lambda h+\mu h]+[\lambda'h+\mu h]+[\lambda''h+\mu h]+\text{etc.}\]
\[+ [\lambda hh+\mu hh]+\text{etc. etc.}\]
quae expressio omnino $ff$ radices continet.\ Quodsi hie singulae columnae
verticales seorsim in summam colliguntur, manifeste prodit
\[(f,\lambda+\mu)+(f,\lambda'+\mu)+(f,\lambda''+\mu)+\text{etc.}\]
quam expressionem cum $W$ convenire nullo negotio perspicitur, quum
numeri $\lambda$, $\lambda'$, $\lambda''$ etc.\ per hyp.\ ipsis $\lambda$, $\lambda h$, $\lambda hh\dots\lambda h^{f-1}$
secundum modulum $n$ congrui esse debeant (quonam ordine hie nihil
interest) adeoque etiam $\lambda+\mu$, $\lambda'+\mu$, $\lambda''+\mu$ etc.\ ipsis
$\lambda+\mu$, $\lambda h+\mu$, $\lambda hh+\mu\dots\lambda h^{f-1}+\mu$.\ Q.\ E.\ D.

Huic theoremati adiungimus corollaria sequentia:

\textbf{I.}\ Designante $k$ integrum quemcunque, productum ex
$(f,k\lambda)\cdot(f,k\mu)$ erit
\[=(f,k(\lambda+\mu))+(f,k(\lambda'+\mu))+(f,k(\lambda''+\mu))+\text{etc.}\]

\textbf{II.}\ Quum singulae partes, e quibus $W$ constat, vel cum aggregato
$(f,0)$, quod est $=f$, vel cum aliquo ex his $(f,1)$, $(f,g)$, $(f,gg)\dots(f,g^{n-2})$
conveniant, $W$ ad formam sequentem reduci poterit:
\[W = af+b(f,g)+b'(f,gg)+b''(f,g^3)+\dots+b^e(f,g^{n-2})\]
ubi coefficientes $a$, $b$, $b'$ etc.\ sunt integri affirmativi (quorum aliqui
etiam $=0$ esse possunt).\ E.\ g.\ pro $n=19$, $f=6$, $\lambda=1$, $\mu=2$,
fit productum ex $(6,1)\cdot(6,2) = (6,3)+(6,1)+2(6,4) = (6,1)+2(6,2)+2(6,4)$.

\textbf{III.}\ Hinc patet, si productum ex $(f,k\lambda)\cdot(f,k\mu)$ in formam
similem redigatur atque in $(f,\lambda)\cdot(f,\mu)$, coefficientes omnino eosdem
prodire debere, i.\ e.\ productum illud fieri
\[af+b(f,k)+b'(f,kg)+b''(f,kgg)\dots\]

\textbf{IV.}\ Denique perspicuum est, functionem $W$, respectu omnium
radicum in $Q$ contentarum, invariabilem esse, quum hae omnes in
periodis $(f,\lambda+\mu)$, $(f,\lambda'+\mu)$ etc.\ contentae sint, quae ad summam
$W$ conspirant, et quae omnes, quum $f$ periodos similes componentes, vel
inter se diversae sint, vel saltem partes similes prodeant, quarum unaquaeque
factorem $t^2-ut\cdot v+v^2$ implicat.

\subsubsection*{346.}

\textbf{Theorema.}\ Si perinde ut in art.\ 345 multiplicetur $(f,\lambda)$ in
$(f,\mu)$, sed ita ut alter factor multiplicandus ex singulis radicibus
alterius periodi similis $(f,\nu)$ constet, puta sit
\[(f,\mu) = (f,\nu)+(f,\nu g^f)+(f,\nu g^{2f})\dots+(f,\nu g^{(e-1)f})\]
productum erit aggregatum $e$ talium aggregationum qualium $W$ est in
art.\ 345, puta
\[W' = W+(\nu g^f)+(W+(\nu g^{2f}))+\text{etc.}\]
u bi $(W+(\nu g^f))$ etc.\ exprimit aggregationem, quae oritur ex $W$, si pro
singulis radicibus $[\lambda]$, $[\lambda']$ etc.\ substituitur $[\lambda g^f]$, $[\lambda' g^f]$ etc.
adeoque pro singulis periodis $(f,\lambda+\mu)$ substituitur $(f,\lambda g^f+\mu)$ etc.

Huius theorema demonstratio, quae ex art.\ 345 nullo negotio derivatur,
supersedere possumus.

\subsubsection*{347.}

Sit $F$ functio invariabilis indeterminatarum $t$, $u$, $v$ etc., puta ea,
quae non mutatur, quomodocunque indeterminatae inter se permutentur;
cuiusmodi sunt e.\ g.\ summa omnium, summa productorum e binis,
summa productorum e ternis etc.\ *) Sit porro $W$ aggregatum $f$
periodorum similium, puta
\[W = (f,\lambda)+(f,\lambda')+(f,\lambda'')+\text{etc.}\]
et substituantur aggregata $(f,\lambda)$, $(f,\lambda')$, $(f,\lambda'')$ etc.\ pro
indeterminatis $t$, $u$, $v$ etc.\ resp., eius valorem per praecepta art.\ 345
reduci ad
\[W = A+a(f,1)+a'(f,g)+a''(f,gg)+\dots+a^e(f,g^{n-2})\]
tum dico, si $F$ sit functio invariabilis, eas, quae sub eadem periodo
$ef$ terminorum contentae sunt, e.\ g.\ generatim tales $(f,\lambda g^{\mu e})$ et
$(f,\lambda g^{\nu e+\sigma})$, designante $\sigma$ integrum quemcunque, eosdem habituras
esse coefficientes.

\textbf{Dem.}\ Quum periodus $(ef,\lambda g^{\mu e})$ cum hac $(ef,\lambda)$ sit identica,
hae pro $f$, $u$, $V$ etc.\ resp.\ substitutis, idem prorsus producent quam
illis, e quibus posterior constat, etsi alio ordine.\ Hinc manifestum est,
$F$ in $W$ transire supponitur, coefficientes periodorum $(f,\lambda g^{\mu e})$,
$(f,\lambda' g^{\mu e})$ etc.\ cum iis $(f,\lambda g^{\nu e})$, $(f,\lambda' g^{\nu e})$ etc.\ quae in $W$
coincidunt.\ At per art.\ 345 erit
\[W = A+a(f,g^\mu)+a'(f,g^{\mu+1})+a''(f,g^{\mu+2})\dots+a^e(f,g^{\mu+e-1})\]
\[= A+a(f,g^\nu)+a'(f,g^{\nu+1})+a''(f,g^{\nu+2})\dots+a^e(f,g^{\nu+e-1})\]
quare quum haec expressio cum $W$ convenire debeat, coefficientes periodorum
$(f,g^{\mu})$, $(f,g^{\mu+1})$ etc.\ (incipiendo ab $a$) cum $(f,g^{\nu})$,
$(f,g^{\nu+1})$ etc.\ necessario conveniet, unde nullo negotio concluditur,
generaliter coefficientes periodorum $(f,\lambda g^{\mu e})$, $(f,\lambda' g^{\mu e})$ etc.,
quae sunt $(\mu+\lambda e)$, $(\mu+\lambda' e)$ etc.\ cum $(f,\lambda g^{\nu e})$,
$(f,\lambda' g^{\nu e})$ etc., quae sunt $(\nu+\lambda e)$, $(\nu+\lambda' e)$ etc.\ inter se
convenire debere.\ Q.\ E.\ D.

Hinc manifestum est, $W$ reduci posse ad formam
\[W = A+a(ef,1)+a'(ef,g)+a''(ef,gg)\dots+a^e(ef,g^{e-1})\]
um bi omnes coefficientes $A$, $a$ etc.\ integri erunt, si omnes coefficientes
determinati in $F$ sunt integri.

Porro facile perspicietur, si postea pro indeterminatis in $F$ substituantur
$e$ periodi $f$ terminorum in alia periodo $ef$ terminorum, puta in
$(ef,\lambda)$ contentae, quae manifesto erunt $(f,\lambda)$, $(f,\lambda g^e)$,
$(f,\lambda g^{2e})\dots(f,\lambda g^{(e-1)e})$, inde prodeuntem fore
\[A+a(ef,\lambda)+a'(ef,\lambda g)+a''(ef,\lambda gg)\dots+a^e(ef,\lambda g^{e-1})\]

\subsubsection*{348.}

Quum in aequatione quacunque
\[x^f-\alpha x^{f-1}+\beta x^{f-2}-\gamma x^{f-3}+\text{etc.} = 0\]
coefficientes $\alpha$, $\beta$, $\gamma$ etc.\ sint functiones invariabiles radicum, puta
$\alpha$ summa omnium, $\beta$ summa productorum e binis, $\gamma$ summa
productorum e ternis etc., in aequatione, cuius radices sunt radices in
periodo $(f,\lambda)$ contentae, coefficiens primus erit $=(f,\lambda)$, singuli
reliqui vero sub formam talem reduci poterunt, ubi omnes $A$, $a$,
$a'$ etc.\ erunt integri; praetereaque patet, aequationem, cuius radices sint
radices in quacunque alia periodo $(f,k\lambda)$ contentae, ex illa derivari, si
in singulis coefficientibus pro $(f,1)$, $(f,g)$, $(f,gg)$ etc.\ substituantur
$(f,k)$, $(f,kg)$, $(f,kgg)$ etc.

Hoc itaque modo assignari poterunt aequationes, quarum radices sint
radices contentae in $(f,1)$, $(f,g)$, $(f,gg)$ etc., quamprimum e aggregata
$(f,1)$, $(f,g)$, $(f,gg)$ etc.\ innotuerunt, aut potius quamprimum unum
quodcunque eorum inventum est, quoniam per art.\ 346 ex uno omnia
reliqua rationaliter deducere licet.\ Quo pacto simul functio $X$ in
factores $f$ dimensionum resoluta habetur, productum enim e functionibus
$z$, $z'$, $z''$ etc.\ manifesto erit $=X$.

\textbf{Ex.}\ Pro $n=19$ summa omnium radicum in periodo $(6,1)$
est $=(6,1)=\alpha$; summa productorum e binis fit $=(6,1)+(6,4)=\beta$;
similiter summa productorum e ternis $=2+(6,2)=\gamma$; summa
productorum e quaternis $=3+(6,1)+(6,4)=\delta$; summa productorum
e quinis $=(6,1)+(6,2)=\varepsilon$; productum ex omnibus $=1$; quare
aequatio $z^6-\alpha z^5+\beta z^4-\gamma z^3+\delta zz-\varepsilon z+1=0$
omnes radices in $(6,1)$ contentas complectitur.

Quodsi in coefficientibus pro $(6,1)$, $(6,2)$, $(6,4)$ resp.\ substituantur
$(6,2)$, $(6,4)$, $(6,1)$, prodibit aequatio $z'=0$, quae radices in $(6,2)$
complectetur; et si eadem commutatio hic denuo applicatur, habebitur
aequatio $z''=0$, radices in $(6,4)$ complectens, productumque $zz'z''$
erit $=X$.

\subsubsection*{349.}

Plerumque commodius est, praesertim quoties $f$ est numerus magnus,
coefficientes $\beta$, $\gamma$ etc.\ secundum theorema \textsc{Newtonianum} e summis
potestatum radicum in $(f,\lambda)$ cura deducere.\ Scilicet sponte patet,
summam quadratorum radicum in $(f,\lambda)$ contentarum esse $(f,2\lambda)$,
summam cuborum $(f,3\lambda)$, summam biquadratorum $(f,4\lambda)$ etc.\ Scribendo
itaque brevitatis caussa pro $(f,\lambda)$, $(f,2\lambda)$, $(f,3\lambda)$ etc.\ $q$, $q'$,
$q''$ etc.\ erit
\[\alpha=q,\quad 2\beta=\alpha q-q',\quad 3\gamma=\beta q-\alpha q'+q'',\quad\text{etc.}\]
um bi producta e duabus periodis per art.\ 345 statim in summas
convertenda.\ Ita in exemplo nostro, scribendo pro $(6,1)$, $(6,2)$,
$(6,4)$ resp.\ $p$, $p'$, $p''$, fit
\begin{align*}
q &= p,\quad q'=p',\quad q''=p'',\quad q'''=p,\quad q^{iv}=p',\quad q^v=p'' \\
2\beta &= pp'-p'' = 6+2p+p' \\
4\gamma &= (2+2p+p')p-(2+p+p')p'+pp'' = 12+4p+4p'' \\
48 &= (2+2p+p')p-(2+p+p')p'+pp'' = \text{etc.}
\end{align*}
Ceterum sufficit semissem coefficientium tantum hoc modo computare;
etenim non difficile probatur, ultimos ordine inverso primis vel aequales
esse, puta ultimum $=1$, penultimum $=\alpha$, antepenultimum $=\beta$
etc., vel ex iisdem resp.\ deduci, si pro $(f,1)$, $(f,g)$ etc.\ substituantur
$(f,-1)$, $(f,-g)$ etc., sive $(f,n-1)$, $(f,n-g)$ etc.

Casus prior locum habet, quando $f$ est par; posterior, quando $f$
impar; coefficiens ultimus autem semper fit $=1$.

Fundamentum huius rei innititur theoremati art.\ 79; sed brevitatis
caussa huic argumento non immoramur.

\subsubsection*{350.}

\textbf{Theorema.}\ Retentis omnibus signis art.\ praec.\ $n-1$ productum
$e$ tribus integris positivis $a$, $b$, $c$, quae est periodus $(abc,\lambda)$,
constet e $\beta\gamma$ periodis $\gamma$ terminorum, puta in $(bc,\lambda)$,
$(bc,\lambda g^a)$, $(bc,\lambda g^{2a})\dots$; sit porro $F=\varphi(t,u,v\dots)$ functio
indeterminatarum $t$, $u$, $v$ etc., similiter affecta ut in art.\ 347, puta in
substituantur aggregata $(bc,\lambda)$, $(bc,\lambda g^a)$, $(bc,\lambda g^{2a})\dots$
pro indeterminatis $t$, $u$, $v$ etc.\ resp., eius valorem per praecepta art.\ 345
et IV reduci ad
\[W = A+a(bc,1)+a'(bc,g)+a''(bc,gg)\dots+a^e(bc,g^{ab-1})\]
Tum dico, si $F$ sit functio invariabilis, eas periodos in $W$, quae sub eadem
periodo $abc$ terminorum contentae sunt, e.\ g.\ generatim tales
$(bc,\lambda g^{\mu a})$ et $(bc,\lambda g^{\nu a+\sigma})$, designante $\sigma$ integrum quemcunque,
eosdem habituras esse coefficientes.

\textbf{Dem.}\ Quum periodus $(abc,\lambda g^{\mu a})$ cum hac $(abc,\lambda)$ sit
identica, hae pro $f$, $u$, $V$ etc.\ resp.\ substitutis, idem prorsus
producent quam illis, e quibus posterior constat, etsi alio ordine.\ Unde
manifestum est, $W$ in $F$ transire supponitur, coefficientes periodorum
$(bc,\lambda g^{\mu a})$, $(bc,\lambda' g^{\mu a})$ etc.\ cum iis $(bc,\lambda g^{\nu a})$,
$(bc,\lambda' g^{\nu a})$ etc.\ quae in $W$ coincidunt.\ At per art.\ 347 erit
\[W = A+a(bc,g^\mu)+a'(bc,g^{\mu+1})+a''(bc,g^{\mu+2})\dots+a^e(bc,g^{\mu+ab-1})\]
\[= A+a(bc,g^\nu)+a'(bc,g^{\nu+1})+a''(bc,g^{\nu+2})\dots+a^e(bc,g^{\nu+ab-1})\]
quare quum haec expressio cum $W$ convenire debeat, coefficientes
primus, secundus, tertius etc.\ (incipiendo ab $a$) cum
$a+1^{ta}$, $a+2^{ta}$, $a+3^{ta}$ etc.\ necesario conveniet, unde nullo negotio
concluditur, generaliter coefficientes periodorum
$(bc,\lambda g^{\mu a+\sigma})$, $(bc,\lambda' g^{\mu a+\sigma})$, $(bc,\lambda'' g^{\mu a+\sigma})$ etc.,
quae sunt $(\mu+\lambda a+\sigma)$, $(\mu+\lambda' a+\sigma)$ etc.\ inter se
convenire debere.\ Q.\ E.\ D.

Hinc manifestum est, $W$ reduci posse ad formam
\[W = A+a(abc,1)+a'(abc,g)+a''(abc,gg)\dots+a^e(abc,g^{c-1})\]
um bi omnes coefficientes $A$, $a$ etc.\ integri erunt, si omnes coefficientes
determinati in $F$ sunt integri.

Porro facile perspicietur, si postea pro indeterminatis in $F$ substituantur
$b$ periodi $c$ terminorum in alia periodo $abc$ terminorum, puta in
$(abc,\lambda k)$ contentae, quae manifesto erunt $(c,\lambda k)$,
$(c,\lambda kg^a)$, $(c,\lambda kgg^a)\dots(c,\lambda kg^{(b-1)a})$, inde prodeuntem fore
\[A+a(abc,k)+a'(abc,kg)+a''(abc,kgg)\dots+a^e(abc,kg^{c-1})\]

\subsubsection*{351.}

Retentis itaque omnibus signis art.\ praec., manifestum est, singulos
coefficientes aequationis, cuius radices sunt aggregata $(bc,\lambda)$,
$(bc,\lambda g^a)$, $(bc,\lambda g^{2a})\dots$, sub formam $A+a(abc,1)+a'(abc,g)\dots$
reduci posse, atque numeros $A$, $a$ etc.\ omnes fieri integros; aequationem
autem, cuius radices sint $b$ periodi $c$ terminorum in aliqua periodo data
$abc$ terminorum contentae, ex illa derivari, si ubique in coefficientibus pro
qualibet periodo $(abc,\mu)$ substituatur $(abc,\lambda\mu)$.

Si igitur $a>1$, omnes $b$ periodi $c$ terminorum determinabuntur
per aequationem $b^{ti}$ gradus, cuius singuli coefficientes sub formam
$A+a(abc,1)\dots$ rediguntur, adeoque sunt quantitates cognitae, quoniam
$(abc,1)=-(abc,g)=-1$.

Si vero $a>1$, $b>1$, coefficientes aequationis, cuius radices sunt
omnes periodi $c$ terminorum in aliqua periodo data $abc$ terminorum
contentae, quantitates cognitae erunt, simulac valores numerici omnium
$a$ periodorum $bc$ terminorum innotuerunt.

Ceterum calculus coefficientium harum aequationum saepe commodius
instituitur, praesertim quando $c$ non est valde parvus, si primo summae
potestatum radicum eruuntur, ac dein ex his per theorema \textsc{Newtonianum}
coefficientes deducuntur, simili modo ut supra art.\ 349.

\textbf{Ex.\ I.}\ Quaeritur pro $n=19$ aequatio, cuius radices sint aggregata
$(6,1)$, $(6,2)$, $(6,4)$.\ Designando has radices per $p$, $p'$, $p''$
resp., et aequationem quaesitam per $x^3-Axx+Bx-C=0$, fit
$A=p+p'+p''$, $B=pp'+pp''+p'p''$, $C=pp'p''$.\ Hinc $A=(18,1)=-1$;
porro habetur
\begin{align*}
pp &= p+2p'+3p'' \quad pp'=2+p'+2p \quad pp''=2+p+2p' \\
pp'p'' &= (2+p+2p')p'' = 3(6,0)+(6,1)+(6,2)+(6,4) = 18-11 = 7
\end{align*}
um de $B = pp'+pp''+p'p'' = 6(p+p'+p'') = 6(18,1) = -6$; denique
fit $C=7$.\ Quare aequatio quaesita $x^3+xx-6x-7=0$.

Utendo methodo altera habemus
\[p+p'+p''=-1,\quad pp=p+2p'+3p'',\quad pp'=2+p'+2p,\quad pp''=2+p+2p'\]
um de
\[\frac{1}{2}(p^2+p'^2+p''^2) = -1+p+p'+p'' = -1+(6,1) = -1+p\]
similiterque $p^3+p'^3+p''^3 = 36+34(p+p'+p'')-2(pp'+pp''+p'p'')$ etc.
hinc per theorema \textsc{Newtonianum} eadem aequatio derivatur ut ante.

\textbf{Ex.\ II.}\ Quaeritur pro $n=19$ aequatio, cuius radices sint aggregata
$(2,1)$, $(2,7)$, $(2,8)$.\ Quibus resp.\ per $q$, $q'$, $q''$ designatis,
invenitur
\[q+q'+q''=(6,1),\ qq'+qq''+q'q''=(6,1)+(6,4),\ qq'q''=2+(6,2)\]
um de, retentis signis ex.\ praec., aequatio quaesita erit
$x^3-pxx+(p+p'')x-2-p'=0$.

Aequatio, cuius radices sunt aggregata $(2,2)$, $(2,3)$, $(2,5)$, sub
$(6,2)$ contenta, e praecedente deducitur, substituendo pro $p,p',p''$
resp.\ $p'$, $p''$, $p$, eademque substitutione iterum facta, prodit
aequatio, cuius radices sunt aggregata $(2,4)$, $(2,6)$, $(2,9)$ sub
$(6,4)$ contenta.

\subsection*{Disquisitionibus praec.\ superstruitur solutio aequationis $X=0$.}

\subsubsection*{352.}

Theoremata praecedentia cum consectariis annexis praecipua totius
theoriae momenta continent, modusque valores radicum $Q$ inveniendi
paucis complectemur.

Quum in aequatione $X=0$ radices $Q$ in periodos $f$ terminorum
distribuere liceat, ubi $f$ est divisor ipsius $n-1$, prodeunt primo
aequationes, quarum radices sunt hae periodi ipsae; quarum inventio ad
aequationes gradus inferioris reducitur.\ E.\ g.\ pro $n=19$, ubi
$n-1=3\cdot 2$, primo prodit aequatio cubica, cuius radices sunt tres
periodi $(6,1)$, $(6,2)$, $(6,4)$; deinde pro singulis his periodis aequatio
quadratica, cuius radices sunt aggregata binorum terminorum.

Hinc valor cuiusvis aggregati $(f,\lambda)$ eruere licebit; quem inventum
singulae radices $[\lambda]$, $[\lambda']$ etc.\ per aequationem quadraticam
$xx-(f,\lambda)x+[\lambda][\lambda']=0$ obtinentur, ubi $[\lambda][\lambda']$ est productum
ex duabus radicibus reciprocis, adeoque $=1$.\ Quum autem in his
aequationibus coefficientes signis radicalibus impliciti sint, ambiguitas
de signis tollenda erit per artificia infra explicanda.

\subsubsection*{353.}

\textbf{Exemplum primum pro $n=19$.}

Quum hic sit $n-1=3\cdot 2$, $Q$ in tres periodos senarum radicum
distribui debet.\ Sit radix primitiva $g=2$, cuius potestates secundum
modulum 19 hos residua minima praebent:

\begin{center}
\begin{tabular}{c|cccccccccccccccccc}
$k$ & 0 & 1 & 2 & 3 & 4 & 5 & 6 & 7 & 8 & 9 & 10 & 11 & 12 & 13 & 14 & 15 & 16 & 17 \\
\hline
$2^k$ & 1 & 2 & 4 & 8 & 16 & 13 & 7 & 14 & 9 & 18 & 17 & 15 & 11 & 3 & 6 & 12 & 5 & 10
\end{tabular}
\end{center}

Periodi sex terminorum hae sunt:
\begin{align*}
(6,1) &= [1]+[8]+[64]+[512]+[4096]+[32768] \\
      &= [1]+[7]+[8]+[11]+[12]+[18] \\
(6,2) &= [2]+[3]+[5]+[14]+[16]+[17] \\
(6,4) &= [4]+[6]+[9]+[10]+[13]+[15]
\end{align*}

Per art.\ 351 invenitur aequatio cubica $x^3+xx-6x-7=0$, cuius
radices $p$, $p'$, $p''$ aequales sunt his $(6,1)$, $(6,2)$, $(6,4)$.

Iam ad valores numericos eruendos, unus radix realis huius aequationis
per methodos notas invenitur $=2,507018644\dots$\ Haec erit $=(6,1)$.
Reliquae radices ex divisione aequationis per $x-(6,1)$ obtinentur,
$-(6,2)=1,285142\dots$, $-(6,4)=-3,792160\dots$

His inventis, ad periodos binorum terminorum descendimus.\ Aequatio,
cuius radices sunt $(2,1)$ et $(2,8)$, fit per art.\ 346:
\[xx-(6,1)x+(6,2)=0\]
um de
\[(2,1) = \frac{1}{2}(6,1)+\sqrt{\frac{1}{4}(6,1)^2-(6,2)} = 1,25\dots\]
\[(2,8) = \frac{1}{2}(6,1)-\sqrt{\frac{1}{4}(6,1)^2-(6,2)}\]

Simili modo invenitur
\[(2,2) = \frac{1}{2}(6,2)+\sqrt{\frac{1}{4}(6,2)^2-(6,4)}\]
\[(2,3) = \frac{1}{2}(6,2)-\sqrt{\frac{1}{4}(6,2)^2-(6,4)}\]

Restat ut ambiguitatem de signis tollamus.\ Pro $(2,1)$ et $(2,8)$,
si $(2,1)$ radix realis maior statuitur, productum $(2,1)\cdot(2,4)=(4,5)$
quantitatem realem positivam esse oportet, unde signum radicale pro
$(2,1)$ positive, pro $(2,8)$ negative accipiendum est.

Denique ad radices ipsas descendimus.\ Aequatio $xx-(2,1)x+1=0$
dat
\[[1] = \frac{1}{2}(2,1)+i\sqrt{1-\frac{1}{4}(2,1)^2},\quad [18] = \frac{1}{2}(2,1)-i\sqrt{1-\frac{1}{4}(2,1)^2}\]
Reliquae radices per potestates ipsius $[1]$ habebuntur, vel per resolutionem
aliarum aequationum quadraticarum.

\subsubsection*{354.}

\textbf{Exemplum secundum pro $n=17$.}

Hic habetur $n-1=2\cdot 2\cdot 2\cdot 2$, quamobrem calculus radicum $Q$
ad quatuor aequationes quadraticas reducendus erit.

Pro radice primitiva hic accipiemus numerum 3, cuius potestates
residua minima sequentia secundum modulum 17 suppeditant:

\begin{center}
\begin{tabular}{c|cccccccccccccccc}
$k$ & 0 & 1 & 2 & 3 & 4 & 5 & 6 & 7 & 8 & 9 & 10 & 11 & 12 & 13 & 14 & 15 \\
\hline
$3^k$ & 1 & 3 & 9 & 10 & 13 & 5 & 15 & 11 & 16 & 14 & 8 & 7 & 4 & 12 & 2 & 6
\end{tabular}
\end{center}

Hinc emergunt distributiones sequentes complexus $Q$ in periodos duas
octonorum, quatuor quaternorum, octo binorum terminorum:

Periodi octo terminorum: $(8,1)$ et $(8,3)$.

Periodi quatuor terminorum, sub $(8,1)$ contentae: $(4,1)$, $(4,9)$,
$(4,13)$, $(4,15)$; sub $(8,3)$ contentae: $(4,3)$, $(4,10)$, $(4,5)$,
$(4,11)$.

Periodi duorum terminorum: $(2,1)$, $(2,9)$, $(2,13)$, $(2,15)$, sub
$(8,1)$; $(2,3)$, $(2,10)$, $(2,5)$, $(2,11)$, sub $(8,3)$.

Iam primum quaeritur aequatio, cuius radices sint $(8,1)$, $(8,3)$.\ Per
art.\ 351 fit haec $xx+(8,1)x-4=0$, quum sit $(8,1)+(8,3)=-1$,
productumque $=-4$.

Hinc $(8,1) = -\frac{1}{2}+\frac{1}{2}\sqrt{17}$, $(8,3)=-\frac{1}{2}-\frac{1}{2}\sqrt{17}$.

Determinandum restat, utram radicem per $(8,1)$, utram per $(8,3)$
exprimere oporteat.\ Computando productum $(8,1)\cdot(2,1)$, invenitur
$=(2,1)+(2,8)+(2,9)+(2,13)+(2,15)+(2,16)$, quod cum sub
$(8,1)$ contineatur, quantitatem realem affirmativam esse oportet.\ Porro
$(2,1)=2\cos\frac{2\pi}{17}$, etc., unde computando invenitur $(8,1)$ affirmativum
esse debere, quum sit maior quam $(8,3)$; ergo
\[(8,1) = \frac{-1+\sqrt{17}}{2},\quad (8,3) = \frac{-1-\sqrt{17}}{2}.\]

Secundo, aggregata $(4,1)$, $(4,9)$, $(4,13)$, $(4,15)$ determinantur
per aequationem quarti gradus, cuius radices sunt illa aggregata, quae per
art.\ 351 fit
\[x^4-(8,1)x^3+\tfrac{1}{2}(5+(8,1)+4(8,3))xx-\tfrac{1}{2}(1+3(8,1))x-1=0\]
Facilius tamen haec aggregata per resolutionem duarum aequationum
quadraticarum obtinentur.\ Per art.\ 346 fit productum
$(4,1)\cdot(4,9)=(4,10)+(4,3)+(4,11)+(4,5)=(8,3)$.

Porro $(4,1)+(4,9)=(8,1)$.\ Quare $(4,1)$ et $(4,9)$ sunt radices
aequationis $xx-(8,1)x+(8,3)=0$, unde
\[(4,1) = \tfrac{1}{2}(8,1)+\tfrac{1}{2}\sqrt{(4+(8,1)-2(8,3))}\]
\[(4,9) = \tfrac{1}{2}(8,1)-\tfrac{1}{2}\sqrt{(4+(8,1)-2(8,3))}\]

Eadem methodus etiam hanc formulam largitur $(4,10)=-\frac{1}{2}(8,3)
+\frac{1}{2}\sqrt{(4+(8,3)-2(8,1))}$, etc.

Dubium vero, utram radicem per $(4,1)$, utram per $(4,9)$ exprimere
oporteat, per artificium sequens, cuius mentionem in art.\ 352 iniecimus,
tolletur.

Evolvatur productum ex $(4,1)-(4,9)$ in $(4,3)-(4,10)$, unde
emergere invenietur $=(4,1)-(4,9)\cdot(4,3)-(4,10)=-1+2\sqrt{17}$,
praetereaque etiam producti factor primus $(4,1)-(4,9)$ positivus est,
puta $=+2\sqrt{(4+(8,1)-2(8,3))}$; etiam alter factor
$(4,3)-(4,10)$ positivus esse debet, puta $=+\sqrt{(4+3(8,1)+4(8,3))}$,
quare necessario $(4,1)$ radici priori, in qua signum positivum radicali
praefigitur, et $(4,9)$ posteriori aequale statui.

Ceterum hinc iidem valores numerici derivantur ut supra.

Cunctis aggregatis quatuor terminorum inventis, progredimur ad
aggregata duorum terminorum.

Aequatio (C), cuius radices sunt hae $(2,1)$, $(2,13)$ sub $(4,1)$
contentae, eruitur haec $xx-(4,1)x+(4,9)=0$, sive
\[\tfrac{1}{2}(4,1)\pm\tfrac{1}{2}\sqrt{(-4(4,9)+(4,1)^2)}\]
cuius radices sunt $(2,1)$, $(2,13)$; sub $(2,1)$, $(2,9)$ etc.

Statuimus eam, ubi quantitas radicalis positive sumitur et cuius valor
reperitur $=1,8649444588$, cuius valor $=0,1845367189$.

$=(2,1)$, unde $(2,13)$ aequale fiet alteri.

Si aggregata reliqua duorum terminorum $(2,2)$, $(2,3)$, $(2,4)$,
$(2,5)$, $(2,6)$, $(2,7)$, $(2,8)$ placet per methodum art.\ 346
investigare, eaedem formulae adhiberi poterunt, quae in ex.\ praec.\ pro
quantitatibus similiter designatis tradidimus, puta $(2,2)$, (sive $(2,15)$),
$=(2,1)^2-2$ etc.

Si vero magis arridet, binas per resolutionem aequationis quadraticae
computare, pro his $(2,9)$, $(2,15)$ producitur aequatio
$xx-(4,9)x+(4,10)=0$, cuius radices evolvuntur
$=\frac{1}{2}(4,9)\pm\frac{1}{2}\sqrt{(4+(4,1)-2(4,10))}$; quo pacto vero signum
ambiguum hie definire oporteat, simili modo decidetur ut supra.\ Scilicet
per evolutionem producti $(2,1)-(2,13)$ in $(2,9)-(2,15)$ prodit
$(4,1)+(4,9)-(4,10)-(4,3)$, quod quum manifesto sit negativum,
factor $(2,1)-(2,13)$ vero positivus, necessario $(2,9)-(2,15)$
negativus esse debebit, quocirca in expressione ante data signum superius
positivum pro $(2,15)$, pro $(2,9)$ inferius negativum adoptandum erit.

Hinc computatur $(2,9)=1,4780178344$, $(2,15)=-1,9659461994$,
quum ex Perinde $(2,1)-(2,13)$ in $(2,3)-(2,5)$ prodeat
$(4,9)-(4,10)$, adeoque quantitas positiva, factorem $(2,3)-(2,5)$
positivum esse concludimus; hinc simili calculo ut ante instituto invenitur
\begin{align*}
(2,3) &= \tfrac{1}{2}(4,3)+\tfrac{1}{2}\sqrt{(4+(4,10)-2(4,9))} = 0,8914767116 \\
(2,5) &= \tfrac{1}{2}(4,3)-\tfrac{1}{2}\sqrt{(4+(4,10)-2(4,9))} = -0,5473259801
\end{align*}

Denique per operationes omnino analogas eruitur
\begin{align*}
(2,10) &= \tfrac{1}{2}(4,10)+\tfrac{1}{2}\sqrt{(4+(4,3)-2(4,1))} = -1,7004342715 \\
(2,11) &= \tfrac{1}{2}(4,10)-\tfrac{1}{2}\sqrt{(4+(4,3)-2(4,1))} = -1,2052692728
\end{align*}

Superest ut ad radices $Q$ ipsas descendamus.

Aequatio (D), cuius radices sunt $[1]$ et $[16]$, prodit
$xx-(2,1)x+1=0$, unde radices $\frac{1}{2}(2,1)\pm\frac{1}{2}i\sqrt{(4-(2,1)^2)}$
sive $\frac{1}{2}(2,1)\pm\frac{1}{2}i\sqrt{((2,1)-(2,15))}$; signum superius pro $[1]$,
inferius pro $[16]$ adoptamus.\ Quatuordecim reliquae radices vel per
potestates ipsius $[1]$ habebuntur; vel per resolutionem septem aequationum
quadraticarum, quae singulae binas exhibent, ubi incertitudo de signis
quantitatum radicalium per idem artificium tolli poterit ut in praecedentibus.

Radices $[4]$ et $[13]$ sunt radices aequationis $xx-(2,12)x+1=0$,
adeoque quantitas realis negativa, quare quum $[1]-[16]$ in $[4]-[13]$
autem prodit productum ex imaginaria $i$ in realem positivam, etiam
$[1]-[16]$ in $[4]-[13]$ productum ex $i$ in realem positivam propter
$ii=-1$; hinc colligitur, pro $[4]$ signum superius, pro $[13]$ inferius
accipiendum esse.

Simili modo pro radicibus $[8]$ et $[9]$ invenitur perius, pro $[8]$
signum superius, pro $[9]$ inferius accipere oportet.

Computando perinde radices reliquas, sequentes valores numericos
obtinemus, ubi radicibus prioribus signa superiora, posterioribus inferiora
respondere subintelligendum est:

\begin{center}
\begin{tabular}{r@{ $=$ }l}
$[1]$,$[16]$ & $0,9324722294\pm 0,3612416662i$ \\
$[2]$,$[15]$ & $0,7390089172\pm 0,6736956436i$ \\
$[3]$,$[14]$ & $0,4457383558\pm 0,8951632914i$ \\
$[4]$,$[13]$ & $0,0922683595\pm 0,9957341763i$ \\
$[5]$,$[12]$ & $-0,273662990\pm 0,9618256432i$ \\
$[6]$,$[11]$ & $-0,6026346364\pm 0,7980172273i$ \\
$[7]$,$[10]$ & $-0,8502171357\pm 0,5264321629i$ \\
$[8]$,$[9]$  & $-0,9829730997\pm 0,1837495178i$
\end{tabular}
\end{center}

\subsubsection*{355.}

\textbf{Aggregata, in quibus terminorum multitudo par est, sunt quantitates
reales.}

Quum $n$ semper supponatur impar, erit 2 inter factores ipsius $n-1$;
complexusque $Q$ ex $\frac{1}{2}(n-1)$ periodis duorum terminorum formatus.
Talis periodus ut $(2,\lambda)$, constabit e radicibus $[\lambda]$ et
$[\lambda g^{\frac{1}{2}(n-1)}]$, denotante $g$ ut supra radicem primitivam quamcunque
pro modulo $n$.\ Sed fit $g^{\frac{1}{2}(n-1)}\equiv -1$ (mod.\ $n$) (V.\ art.\ 62),
um de $[\lambda g^{\frac{1}{2}(n-1)}] = [-\lambda]$.

Quare supponendo $[\lambda] = \cos\frac{\lambda P}{n}+i\sin\frac{\lambda P}{n}$, fit
$[-\lambda] = \cos\frac{\lambda P}{n}-i\sin\frac{\lambda P}{n}$, et proin $(2,\lambda) =
2\cos\frac{\lambda P}{n}$.\ Unde hoc loco hanc tantummodo conclusionem deducimus,
valorem cuiusvis aggregati duorum terminorum esse quantitatem realem.

Quum quaevis periodus, cuius terminorum multitudo par $=2a$, in $a$
periodos binorum terminorum discerpi possit, patet generalius, valorem
cuiusvis aggregati, in quo terminorum multitudo par est, esse quantitatem
realem.

\vspace{1cm}
\begin{center}
\rule{0.5\textwidth}{0.4pt}
\end{center}

\end{document}
